% Standalone version assembled from the modular manuscript.
\documentclass[11pt,a4paper]{article}
\usepackage[T1]{fontenc}
\usepackage[utf8]{inputenc}
\usepackage{lmodern}
\usepackage[margin=1.08in]{geometry}
\usepackage{amsmath,amssymb,amsthm,mathtools}
\usepackage{microtype}
\usepackage{xcolor}
\usepackage{enumitem}
\usepackage{booktabs}
\usepackage{hyperref}
\hypersetup{colorlinks=true,linkcolor=blue!45!black,citecolor=green!35!black,
 urlcolor=blue!55!black,pdftitle={Grauert's ruled surfaces and semipositive metrics: Lipschitz existence and a C1 obstruction},pdfauthor={}}
\numberwithin{equation}{section}
\newtheorem{theorem}{Theorem}[section]
\newtheorem{proposition}[theorem]{Proposition}
\newtheorem{lemma}[theorem]{Lemma}
\newtheorem{corollary}[theorem]{Corollary}
\theoremstyle{definition}
\newtheorem{definition}[theorem]{Definition}
\newtheorem{assumption}[theorem]{Assumption}
\theoremstyle{remark}
\newtheorem{remark}[theorem]{Remark}
\newcommand{\C}{\mathbb C}
\newcommand{\R}{\mathbb R}
\newcommand{\Z}{\mathbb Z}
\renewcommand{\P}{\mathbb P}
\newcommand{\OO}{\mathcal O}
\newcommand{\ddc}{dd^{c}}
\newcommand{\cOne}{c_1}
\newcommand{\wt}{\widetilde}
\DeclareMathOperator{\Hom}{Hom}
\DeclareMathOperator{\Pic}{Pic}
\DeclareMathOperator{\PSH}{PSH}
\DeclareMathOperator{\Null}{Null}
\DeclareMathOperator{\Bs}{Bs}
\DeclareMathOperator{\Supp}{Supp}
\DeclareMathOperator{\Div}{div}
\DeclareMathOperator{\Ext}{Ext}
\DeclareMathOperator{\Rea}{Re}
\DeclareMathOperator{\Ima}{Im}
\setlist{itemsep=3pt,topsep=5pt}
\setlength{\emergencystretch}{2em}
\title{Grauert's ruled surfaces and semipositive metrics:\\
 Lipschitz existence and a \texorpdfstring{$C^1$}{C1} obstruction}
\author{}
\date{3 November 2026}
\begin{document}
\maketitle
\begin{abstract}
Let $B$ be a smooth projective complex curve of genus at least two, let
$F$ be a positive-degree line bundle on $B$, and let $0\ne\xi\in H^1(B,F)$.
The affine $F$-torsor defined by $\xi$ has a projective completion
$q:X=\P_{\mathrm{lines}}(E)\to B$, where
$0\to F\to E\to\OO_B\to0$ represents $\xi$. Its infinity section
$C=\P_{\mathrm{lines}}(F)$ has negative normal bundle. We give a detailed
account of this construction and of the distinguished line bundle
$L_0=\OO_X(C)\otimes q^*F$. We prove that $L_0$ is nef and big, has null
locus $C$, and is not semiample. We then prove two regularity statements:
$L_0$ admits a Hermitian metric with locally Lipschitz weights and positive
curvature current, smooth with strictly positive curvature away from $C$;
but it admits no such semipositive metric of class $C^1$, even near $C$.
The existence proof uses a quantitatively controlled cusp in the norm of a
defining section of $C$ and a regularized maximum. The obstruction uses
holomorphic graph tests for first derivatives of plurisubharmonic weights.
The resulting bounded potentials have minimal singularities and answer,
for both the original and modified examples, the bounded-potential question
recorded in Zhang's December 2025 preprint. All regularity arguments are
proved here. The account distinguishes them from Grauert's contraction
construction, Koike's smooth first-obstruction theorem, and Zhang's
calculation for the original example; it makes no claim about publication
priority of the arguments presented here.
\end{abstract}
\tableofcontents

\section{Introduction and principal results}\label{sec:introduction}
Grauert's construction in \cite[\S4,8(d), pp.~365--366]{Grauert1962}
exhibits a smooth projective ruled surface containing a negative section
whose analytic contraction is not projective. The construction is governed
by an additive cohomology class, not just by the numerical data of the
section. This distinction is visible in the positivity of a particular
line bundle on the ruled surface. Filip and Tosatti singled out this
example in \cite[Problem~2.2, p.~2989]{FilipTosatti2018}. Koike's
first-obstruction theorem and his modified example appear in
\cite[Theorem~1.4 and Example~4.1]{Koike2021}. Zhang treats the original
degree-one construction in \cite[Theorem~3.2]{Zhang2025} and asks in
\cite[\S4]{Zhang2025} whether its distinguished class, or that of the
modified example, has a positive current with bounded potentials and
minimal singularities.

This manuscript works with the following data throughout its geometric
applications. Let $B$ be a smooth connected projective complex curve of
genus $g\ge2$, let $F$ be a line bundle of degree $d>0$, and choose
$0\ne\xi\in H^1(B,F)$. Let
\begin{equation}\label{eq:main-extension}
0\longrightarrow F\longrightarrow E\longrightarrow\OO_B\longrightarrow0
\end{equation}
be the extension represented by $\xi$, and put
\begin{equation}\label{eq:main-data}
 X=\P_{\mathrm{lines}}(E),\qquad C=\P_{\mathrm{lines}}(F),\qquad
 L_0=\OO_X(C)\otimes q^*F.
\end{equation}
Here projectivization parametrizes lines in each vector space $E_b$.
The notation $F,E$ is used for vector bundles and for their sheaves of
holomorphic sections; the affine torsor itself is denoted by $T$.
The construction and every transition convention are specified in
Section~\ref{sec:geometry}.

\begin{theorem}\label{main:regularity}
For the data in \eqref{eq:main-extension}--\eqref{eq:main-data} the following
statements hold.
\begin{enumerate}[label=\textup{(\roman*)}]
\item The surface $X$ is smooth and projective, $C\cong B$, and
$N_{C/X}\cong F^{-1}$. In particular $C^2=-d$.
\item The bundle $L_0$ is nef and big,
\[
L_0^2=d,\qquad L_0\cdot C=0,\qquad
L_0|_C\cong\OO_C,\qquad \Null(L_0)=C.
\]
Every section of every positive tensor power of $L_0$ vanishes along $C$;
in particular $L_0$ is not semiample.
\item There exists a positive definite Hermitian metric on $L_0$ with
locally Lipschitz weights and semipositive curvature in the sense of
currents. The metric can be chosen smooth with strictly positive curvature
on $X\setminus C$. Its curvature has bounded local potentials and minimal
singularities.
\item No neighborhood of $C$ carries a $C^1$ Hermitian metric on $L_0$
whose curvature is a positive current. In particular there is no
$C^2$ or smooth semipositive Hermitian metric on $L_0$.
\end{enumerate}
\end{theorem}

The proof is distributed as follows. The geometric statements and the
absence of sections on the affine torsor are proved in
Section~\ref{sec:geometry}. The obstruction in (iv) is
Theorem~\ref{obs:c1-obstruction}. The metric in (iii) is constructed in
Theorem~\ref{met:lipschitz-global}, and its application to $L_0$ is given in
Corollary~\ref{con:application}. This organization makes the existence and
nonexistence arguments independent: neither is deduced by a regularization
or limiting argument from the other.

The contrast in Theorem~\ref{main:regularity} concerns the regularity of local
metric weights. It is not a claim that every notion of differentiability
between Lipschitz continuity and $C^1$ has been classified. In particular,
existence of a Lipschitz metric does not supply a semipositive $C^1$
approximation in the same class.

\paragraph{Scope of self-containedness.}
All local analytic lemmas, the curvature inequalities, the first-jet
obstruction, the metric gluing, and the minimal-singularity comparison are
proved in detail. For the surrounding algebraic and analytic geometry we
use standard foundational results: curve Riemann--Roch and Serre duality,
solvability of the Poisson equation on a compact curve, the algebraic--analytic
comparison for projective varieties, the Nakai--Moishezon criterion,
the Hilbert polynomial of an ample line bundle, Grauert's contraction
criterion, the Leray--Hirsch theorem, and the compact K\"ahler
$\partial\bar\partial$-lemma. Each use is stated with its needed hypotheses
and cited. These background theorems are not presented as new results or
silently replaced by an appeal to a stronger assertion. In particular,
the metric existence and obstruction proofs do not invoke a holomorphic
tubular neighborhood theorem or a smooth semipositivity theorem.

% Begin preliminaries.tex
\section{Conventions and elementary analytic tools}\label{sec:preliminaries}

\subsection{Metric weights, curvature, and regularity}
Let $L$ be a holomorphic line bundle on a complex manifold $S$. A positive
definite Hermitian metric $h$ has, in a nowhere vanishing holomorphic frame
$e_i$, a real weight $\varphi_i$ defined by
\[
 |e_i|_h^2=e^{-\varphi_i}.
\]
If $e_i=g_{ij}e_j$, where $g_{ij}$ is holomorphic and nowhere zero, then
\begin{equation}\label{eq:pre-weight-transition}
 \varphi_i=\varphi_j-\log|g_{ij}|^2.
\end{equation}
The second term is smooth and pluriharmonic. Consequently continuity,
local Lipschitz continuity, and membership in $C^k$ of the weights are
independent of the chosen holomorphic frames. For a continuous positive
metric, its local positive norm functions and its weights have the same
local Lipschitz or $C^k$ regularity: on each relatively compact chart the
norm is bounded above and bounded away from zero, and the logarithm and
exponential have bounded derivatives on the corresponding ranges.

We fix the normalization
\[
 d^c=\frac{i}{4\pi}(\bar\partial-\partial),\qquad
 \ddc=\frac{i}{2\pi}\partial\bar\partial.
\]
For a smooth metric, $c_1(L,h)=\ddc\varphi_i$ is a global real closed
$(1,1)$-form representing $c_1(L)$. More generally, a locally integrable
weight defines the curvature current by the same formula. The metric is
\emph{semipositive in the sense of currents} when its local weights are
plurisubharmonic. For continuous weights, this is equivalent to positivity
of their distributional complex Hessians. A singular metric may have
weights with value $-\infty$; a continuous positive definite metric in this
paper always has finite weights.

Fix a smooth reference metric $h_0$ and write
$\theta=c_1(L,h_0)$. A global function $u$ changes the squared metric by
$h=h_0e^{-u}$ and changes the local weights by
$\varphi_i=\varphi_{0,i}+u$. We write
\[
 \PSH(S,\theta)=\{u:\ \varphi_{0,i}+u\text{ is plurisubharmonic
 on every local frame chart}\}.
\]
Equivalently $\theta+\ddc u\ge0$ in currents, with the usual upper
semicontinuous representative. This notation makes sense on an open
subset as well as on a compact manifold.

For matrix computations, if a real $(1,1)$-form is written
\[
 \beta=\frac{i}{2\pi}\sum_{j,k}\beta_{j\bar k}
 dz_j\wedge d\bar z_k,
\]
we denote its associated Hermitian form by
$\beta_{V\bar W}=\sum_{j,k}\beta_{j\bar k}V^j\overline{W^k}$.
This notation removes the factor $i/(2\pi)$ and does not mean the literal
evaluation of the alternating two-form on $(V,\bar W)$. Positivity of
$\beta$ is positivity of this Hermitian form. The frame used to evaluate
it may be smooth rather than holomorphic.

\subsection{Subharmonic limits, maxima, and removable sets}
We record the elementary tools used below. They can also be found in the
general treatments \cite[Chapter~I, \S5]{DemaillyCADG} and
\cite{Demailly2012}.

\begin{lemma}\label{pre:subharmonic-limits}
A locally uniform limit of continuous subharmonic functions is
subharmonic. If both a continuous real function and its negative are
subharmonic, the function is harmonic. A continuous subharmonic function
on a compact connected Riemann surface is constant.
\end{lemma}
\begin{proof}
For a closed Euclidean disc contained in a coordinate domain, the
submean inequality for each approximating function passes to the limit
because convergence is uniform on the closed disc and its boundary.
The limit is continuous and satisfies the submean inequality on every
such disc, which is the characterization of subharmonicity for continuous
functions. If the same inequality holds for a function and its negative,
both inequalities are equalities. The function has the mean-value
property and is harmonic. For completeness, convolving with a radial
smooth approximate identity on a smaller disc leaves a function with
the mean-value property unchanged: the convolution is a weighted average
of its circular means. The convolution is smooth, so the function is
smooth and its Laplacian is zero. Finally, a continuous function on a
compact surface attains its maximum. The maximum principle for the
submean inequality makes a subharmonic function constant on a
neighborhood of a maximum point; the same argument and connectedness
make it constant on the whole surface.
\end{proof}

\begin{lemma}\label{pre:removable}
Let $D$ be a smooth complex hypersurface in a complex manifold and let
$v$ be continuous on a neighborhood of $D$. If $v$ is plurisubharmonic
off $D$, then it is plurisubharmonic everywhere on that neighborhood.
\end{lemma}
\begin{proof}
It suffices to work in a polydisc with coordinates $(z,w)$ in which
$D=\{w=0\}$. A restriction to an affine complex line not contained in
$D$ is subharmonic outside at most one point and is continuous at that
point. Such a point is removable. One direct verification is to take
a small disc around it, let $H$ be the harmonic function with the
continuous boundary values of the restricted function, and apply the
maximum principle on punctured annuli to
$v-H+\varepsilon\log(|\zeta|/R)$. This expression is negative on the
inner boundary for small enough inner radius, and is zero on the outer
boundary. It is therefore nonpositive on the annulus. Let the inner
radius tend to zero, then let $\varepsilon$ decrease to zero. This
proves the harmonic comparison inequality on the full disc, hence
subharmonicity across its center.

For a complex line contained in $D$, translate that line by a sequence
of small nonzero constants in the $w$ coordinate. On each relatively
compact disc the restrictions of $v$ to the translated lines converge
uniformly to its restriction to the original line, by continuity of $v$.
The limit is subharmonic by Lemma~\ref{pre:subharmonic-limits}. Thus the
restriction to every affine complex line is subharmonic. The
line-restriction characterization of a continuous plurisubharmonic
function proves the assertion. This proof also shows that the removable
extension takes the original continuous values on $D$.
\end{proof}

\begin{lemma}\label{pre:regularized-max}
For every $\varepsilon>0$ there exists a smooth convex function
$\mathcal M_\varepsilon:\R^2\to\R$ with the following properties:
\begin{enumerate}[label=\textup{(\roman*)}]
\item it is nondecreasing in each variable;
\item $\mathcal M_\varepsilon(x+a,y+a)=a+\mathcal M_\varepsilon(x,y)$;
\item it equals $\max\{x,y\}$ whenever $|x-y|\ge\varepsilon$;
\item if $u,v$ are continuous plurisubharmonic functions, then
$\mathcal M_\varepsilon(u,v)$ is plurisubharmonic;
\item the maximum and the regularized maximum of two locally Lipschitz
functions are locally Lipschitz.
\end{enumerate}
The same operations preserve $\theta$-plurisubharmonicity when the two
potentials are measured relative to the same smooth form $\theta$.
\end{lemma}
\begin{proof}
Choose an even nonnegative smooth function $\chi$ supported in
$(-\varepsilon,\varepsilon)$ with integral one and put
\[
 \rho_\varepsilon(t)=\int_{\R}|t-a|\chi(a)\,da,
 \qquad
 \mathcal M_\varepsilon(x,y)=\frac{x+y+\rho_\varepsilon(x-y)}2.
\]
Convolution makes $\rho_\varepsilon$ smooth. Convexity of the absolute
value makes it convex, and $|\rho_\varepsilon'|\le1$. When
$|t|\ge\varepsilon$, every $t-a$ has the sign of $t$, and evenness of
$\chi$ gives $\int a\chi(a)\,da=0$. Hence
$\rho_\varepsilon(t)=|t|$ there. These facts prove (i)--(iii): the two
first partial derivatives are
$(1+\rho_\varepsilon')/2$ and $(1-\rho_\varepsilon')/2$, and the Hessian
is a nonnegative multiple of
$\left(\begin{smallmatrix}1&-1\\-1&1\end{smallmatrix}\right)$.

To prove (iv), restrict to a complex disc. The submean inequalities for
$u,v$, monotonicity in each variable, and Jensen's inequality for the
convex function $\mathcal M_\varepsilon$ give the submean inequality for
$\mathcal M_\varepsilon(u,v)$. The function is continuous, so it is
subharmonic on the disc. All disc restrictions have this property, which
proves plurisubharmonicity. The proof for the ordinary maximum follows
either from its submean inequality or by letting $\varepsilon$ tend to
zero locally uniformly.

For (v), the first derivatives of $\mathcal M_\varepsilon$ lie in $[0,1]$
and sum to one. The fundamental theorem of calculus along a segment
therefore gives
\[
 |\mathcal M_\varepsilon(x_1,y_1)-\mathcal M_\varepsilon(x_2,y_2)|
 \le |x_1-x_2|+|y_1-y_2|.
\]
The same inequality holds for the maximum. Composing this inequality
with local Lipschitz bounds for $u,v$ proves (v).

For the last assertion, add a local smooth potential $\varphi_0$ for
$\theta$. Translation invariance gives
\[
 \varphi_0+\mathcal M_\varepsilon(u,v)
 =\mathcal M_\varepsilon(\varphi_0+u,\varphi_0+v).
\]
The right-hand side is plurisubharmonic by (iv). The same identity holds
for the ordinary maximum, proving the assertion on all frame charts.
\end{proof}

\begin{remark}\label{pre:strict-max}
If the two inputs are smooth and strictly plurisubharmonic on a region,
their regularized maximum is smooth and strictly plurisubharmonic there.
Indeed, its complex Hessian is the sum of the two input Hessians with
nonnegative coefficients adding to one, plus a positive semidefinite
quadratic term from the Hessian of $\mathcal M_\varepsilon$. The weighted
sum is positive definite at every point. The same conclusion holds for
strictly $\theta$-plurisubharmonic inputs after adding a local potential
of $\theta$.
\end{remark}

\subsection{Bounded potentials and minimal singularities}
\begin{definition}\label{pre:minimal-definition}
Let $S$ be compact and $\theta$ a smooth closed real $(1,1)$-form.
A function $u\in\PSH(S,\theta)$ has \emph{minimal singularities} if for
every $v\in\PSH(S,\theta)$, $v\not\equiv-\infty$, there exists a constant
$C_v$ such that $v\le u+C_v$ on $S$. The corresponding current
$\theta+\ddc u$ is said to have minimal singularities.
\end{definition}
This is the usual comparison of singularities; see
\cite[Definition~1.4 and Theorem~1.5]{DPS2001} and
\cite[\S1.4]{BEGZ2010}. It does not require that $u$ equal a specially
normalized envelope.

\begin{lemma}\label{pre:bounded-minimal}
Every bounded $u\in\PSH(S,\theta)$ on a compact connected manifold has
minimal singularities. In particular every continuous such $u$ has
minimal singularities and defines a current with bounded local potentials.
\end{lemma}
\begin{proof}
A $\theta$-plurisubharmonic function $v\not\equiv-\infty$ is locally
bounded above: add a local smooth potential of $\theta$ and use the local
upper boundedness of a plurisubharmonic function. A finite relatively
compact coordinate cover of the compact space gives $\sup_Sv<\infty$.
If $u$ is bounded below, then for every point $x$,
\[
 v(x)\le\sup_Sv
 \le u(x)+\sup_Sv-\inf_Su.
\]
This is the defining inequality with $C_v=\sup_Sv-\inf_Su$.
Continuity on the compact space implies boundedness. A local potential
of the current is $\varphi_0+u$; it is bounded on each relatively compact
frame chart because $\varphi_0$ is smooth and $u$ is bounded.
\end{proof}

% End preliminaries.tex
% Begin geometry.tex
\section{The geometry of Grauert's example}\label{geo:section}\label{sec:geometry}

We describe the construction with a positive line bundle of arbitrary
degree.  The original example in \cite[\S4, paragraph 8(d)]{Grauert1962}
uses a curve of genus two and a line bundle of degree one.  Throughout
this section, the projectivization of a vector bundle parametrizes
\emph{lines} in its fibers.

\begin{assumption}\label{geo:data}
Let $B$ be a connected compact Riemann surface of genus $g\geq 2$.
Let $F$ be a holomorphic line bundle on $B$ of degree $d>0$, and choose
a nonzero class
\[
  \xi\in H^1(B,F).
\]
Here and below, $H^1(B,F)$ means cohomology of the sheaf of holomorphic
sections of $F$.
\end{assumption}

The condition $H^1(B,F)\ne 0$ is an additional hypothesis; positivity
of $F$ alone does not imply it.  The following choice provides the
degree-one examples needed here.

\begin{lemma}\label{geo:degree-one-choice}
For any $b_0\in B$, the line bundle $F=\OO_B(b_0)$ has degree one and
\[
  h^1(B,F)=g-1>0.
\]
\end{lemma}

\begin{proof}
A section of $\OO_B(b_0)$ is a meromorphic function on $B$ with no
poles except possibly a simple pole at $b_0$.  If such a function were
nonconstant, it would define a holomorphic map $B\to\P^1$ of degree
one, because the degree equals the total pole order.  A degree-one
map between smooth compact curves is an isomorphism.  This would
contradict $g\geq2$.  Thus $h^0(B,F)=1$.  Riemann--Roch on $B$ now
gives
\[
  h^0(B,F)-h^1(B,F)=\deg F+1-g=2-g,
\]
and the asserted formula follows.
\end{proof}

\subsection{The affine bundle and its projective completion}

Choose a sufficiently fine open cover $\{U_i\}$ of $B$ on which $F$
is trivial, with nowhere-vanishing holomorphic frames $m_i$.  Write
\[
  m_j=g_{ij}m_i
\]
on $U_i\cap U_j$, where the $g_{ij}$ are nowhere-vanishing holomorphic
functions.  Thus fiber coordinates for the total space of $F$ obey
$t_i=g_{ij}t_j$.  Represent $\xi$ by a cocycle
$\{a_{ij}m_i\}$, with $a_{ij}$ holomorphic on $U_i\cap U_j$.  The
cocycle relations are
\begin{equation}\label{geo:cocycle-relations}
  g_{ik}=g_{ij}g_{jk},\qquad
  a_{ik}=a_{ij}+g_{ij}a_{jk}.
\end{equation}

Glue the products $U_i\times\C$ by the affine transformations
\begin{equation}\label{geo:affine-transition}
  t_i=g_{ij}t_j+a_{ij}.
\end{equation}
The result is a holomorphic affine-line bundle $T\to B$.  It is also
called the $F$-torsor associated with $\xi$.  The linear part of its
transition functions is the one defining $F$, while the translation
part records the cohomology class $\xi$.

\begin{lemma}\label{geo:no-section}
The affine bundle $T\to B$ has no holomorphic section.
\end{lemma}

\begin{proof}
A holomorphic section would be represented by holomorphic functions
$s_i$ on $U_i$ satisfying
\[
  s_i=g_{ij}s_j+a_{ij}.
\]
Equivalently,
\[
  a_{ij}m_i=s_im_i-s_jm_j.
\]
This expresses the cocycle representing $\xi$ as a coboundary, contrary
to $\xi\ne0$.  Conversely, a coboundary would give such a section.
\end{proof}

The affine transformations extend to automorphisms of $\P^1$ fixing
the point at infinity.  Adding this point to each fiber gives a smooth
compact surface
\[
  q:X\longrightarrow B.
\]
Its infinity section is denoted by $C$.  Thus $C\cong B$ and
$X\setminus C=T$.  One can identify this completion with a projective
bundle, including its holomorphic extension data, as follows.

\begin{proposition}\label{geo:projective-bundle}
The surface $X$ is the projective bundle $\P_B(E)$ of lines in a
rank-two vector bundle $E$ fitting into the nonsplit exact sequence
\begin{equation}\label{geo:extension}
  0\longrightarrow F\longrightarrow E\longrightarrow\OO_B
    \longrightarrow0.
\end{equation}
Its extension class is $\xi$, up to the harmless sign convention for
identifying extension classes with \v{C}ech cocycles.  The section $C$ is
$\P_B(F)\subset\P_B(E)$.  In particular, $X$ is a smooth projective
surface.
\end{proposition}

\begin{proof}
Define $E$ by identifying vector coordinates according to
\[
  \begin{pmatrix}u_i\\v_i\end{pmatrix}
  =
  \begin{pmatrix}g_{ij}&a_{ij}\\0&1\end{pmatrix}
  \begin{pmatrix}u_j\\v_j\end{pmatrix}.
\]
The relations \eqref{geo:cocycle-relations} imply the matrix cocycle
condition.  The first coordinate axis defines the subbundle $F$,
and the second coordinate gives a quotient map $E\to\OO_B$.
A splitting would provide a lift of $1\in\OO_B$ with coordinates
$(s_i,1)$; its transition rule would be the one excluded in
Lemma~\ref{geo:no-section}.  Hence the extension is nonsplit.

On the chart $v_i\ne0$ of the projectivized bundle, the coordinate
$t_i=u_i/v_i$ satisfies \eqref{geo:affine-transition}.  The complement
of these charts is the section $v_i=0$, namely $\P_B(F)$.  This
identifies the projectivization with the stated completion.

For projectivity, we use two standard algebraic facts explicitly.
A compact Riemann surface is a smooth projective algebraic curve;
and GAGA algebraizes every holomorphic vector bundle on a projective
complex variety \cite{Serre1956}.  Thus $E$ is an algebraic vector
bundle on the projective curve $B$.  Its projective bundle is
projective over $B$, and hence projective over $\C$
\cite{Hartshorne1977}.  Smoothness follows already from its local
description as $U_i\times\P^1$.
\end{proof}

\subsection{The negative section and the distinguished line bundle}

Near $C$, use the fiber coordinate $w_i=1/t_i=v_i/u_i$.  The affine
transition rule becomes
\begin{equation}\label{geo:normal-transition}
  w_i=\frac{w_j}{g_{ij}+a_{ij}w_j},\qquad
  w_j=\frac{g_{ij}w_i}{1-a_{ij}w_i}.
\end{equation}
All these formulas are understood on sufficiently small overlap
neighborhoods of $C$, where the denominators are nonzero.

\begin{proposition}\label{geo:normal-and-L}
The infinity section has normal bundle and self-intersection
\[
  N_{C/X}\cong F^{-1},\qquad C^2=-d.
\]
The distinguished line bundle
\begin{equation}\label{geo:L-definition}
  L_0:=\OO_X(C)\otimes q^*F
\end{equation}
restricts holomorphically trivially to $C$.  With the lines convention
for projectivization, $\OO_X(C)\cong\OO_{\P_B(E)}(1)$.
\end{proposition}

\begin{proof}
Let $I_C$ be the ideal sheaf of $C$.  The class of $w_i$ generates
$I_C/I_C^2$ on $U_i$.  Modulo $I_C^2$, the first formula in
\eqref{geo:normal-transition} reads
\[
  w_i=g_{ij}^{-1}w_j.
\]
This is the same transition rule as $m_i=g_{ij}^{-1}m_j$.
Consequently $I_C/I_C^2\cong F$ under the identification $C\cong B$.
Dualizing gives $N_{C/X}\cong F^{-1}$.  For a smooth divisor in a
smooth surface, its self-intersection equals the degree of its normal
bundle, so $C^2=-d$.

The quotient $E\to\OO_B$ in \eqref{geo:extension}, composed with the
tautological inclusion $\OO_{\P_B(E)}(-1)\to q^*E$, is a section of
$\OO_{\P_B(E)}(1)$.  Its zero divisor is exactly $\P_B(F)=C$, with
multiplicity one.  This proves the assertion about $\OO_X(C)$.
Finally, restriction of a divisor line bundle to its smooth divisor
is the normal bundle.  Therefore
\[
  L_0|_C
   \cong N_{C/X}\otimes F
   \cong F^{-1}\otimes F
   \cong\OO_C.
\]
\end{proof}

For later use, the trivialization on $C$ can be made completely
explicit.  Write $s_C$ for the canonical section of $\OO_X(C)$ and
choose local frames $\sigma_i$ near $C$ with
$s_C=w_i\sigma_i$.  The frames
\[
  \ell_i:=\sigma_i\otimes q^*m_i
\]
of $L_0$ satisfy
\begin{equation}\label{geo:L-transition}
  \frac{\ell_i}{\ell_j}
   =\frac{w_j}{w_i}\,g_{ij}^{-1}
   =1+\frac{a_{ij}}{g_{ij}}w_j
   =\frac{1}{1-a_{ij}w_i}.
\end{equation}
In particular, the frames $\ell_i|_C$ agree on overlaps.  The
translation cocycle disappears upon restriction to $C$ but persists
in the first normal term of the transition functions.

\subsection{Intersection numbers, nefness, and bigness}

We use additive notation for tensor products when computing divisor
classes.  Let $f$ denote a fiber of $q$.  The relevant elementary
intersection identities are
\begin{equation}\label{geo:intersection-table}
  C^2=-d,\quad C\cdot f=1,\quad
  C\cdot q^*F=d,\quad (q^*F)^2=0,\quad q^*F\cdot f=0.
\end{equation}
Indeed, $q|_C$ is an isomorphism, the pullback of a line bundle from
the base has degree zero on a fiber, and two classes pulled back from
a curve have zero intersection on the surface.
More explicitly, $C$ meets each fiber transversely once, and
$q^*F$ is trivial on that fiber. Thus
$L_0|_f\cong\OO_{\P^1}(1)$.

\begin{proposition}\label{geo:nef-big-ample}
The line bundle $L_0$ is nef and big.  Its numerical invariants are
\[
  L_0^2=d,\qquad L_0\cdot C=0,\qquad L_0\cdot f=1.
\]
Moreover, $L_0\cdot D>0$ for every irreducible curve $D\ne C$.
The line bundle
\begin{equation}\label{geo:A-definition}
  A:=L_0^{\otimes2}\otimes\OO_X(-C)
    =\OO_X(C)\otimes q^*F^2
\end{equation}
is ample and satisfies $A^2=3d$ and $A\cdot C=d$.
\end{proposition}

\begin{proof}
The displayed invariants follow from
$L_0=C+q^*F$ and \eqref{geo:intersection-table}.  To check every
curve, let $D\ne C$ be irreducible.  The local intersection
multiplicities of two distinct curves on a smooth surface are
nonnegative, so $C\cdot D\geq0$.

If $q(D)$ is a point, then $D$ is a fiber, since each fiber is an
irreducible copy of $\P^1$.  In this case $L_0\cdot D=1$.
Otherwise, the induced map from the normalization of $D$ to $B$
has a positive degree $n$.  Pullback of a degree-$d$ line bundle has
degree $nd$, and hence
\[
  L_0\cdot D=C\cdot D+nd>0.
\]
Together with $L_0\cdot C=0$, this proves nefness, using the
intersection characterization of nef line bundles on projective
varieties.  It also proves that the null locus, defined here as the
union of curves having zero intersection with $L_0$, is exactly $C$.

For $A=C+2q^*F$, the same calculation gives
\[
  A^2=3d>0,\qquad A\cdot C=d>0,\qquad A\cdot f=1.
\]
For a horizontal irreducible curve $D\ne C$ of degree $n$ over $B$,
one has $A\cdot D=C\cdot D+2nd>0$.  We now apply the
Nakai--Moishezon criterion in its surface form: a line bundle on a
projective surface is ample if its square is positive and its
intersection with every irreducible curve is positive
\cite{Hartshorne1977,Lazarsfeld2004}.  The hypotheses have all been
checked, so $A$ is ample.

For completeness, bigness can be verified directly from growth of
sections, without appealing to the general nef-and-positive-square
criterion.  Multiplication by $s_C^k$ gives an injection
\[
  H^0(X,A^{\otimes k})
    \longrightarrow H^0(X,L_0^{\otimes 2k}),
\]
because $A^{\otimes k}\otimes\OO_X(kC)=L_0^{\otimes2k}$ and $s_C$
does not vanish identically on any open set.  The Hilbert polynomial
and Serre vanishing for an ample line bundle on a projective surface
give
\[
  h^0(X,A^{\otimes k})=\frac{A^2}{2}k^2+O(k)
                       =\frac{3d}{2}k^2+O(k).
\]
Consequently
\[
  \limsup_{m\to\infty}
    \frac{h^0(X,L_0^{\otimes m})}{m^2}\geq\frac{3d}{8}>0.
\]
This is the section-growth definition of a big line bundle.
\end{proof}

\begin{remark}\label{geo:FT-choice}
The formula in \cite[discussion preceding Problem 2.2]{FilipTosatti2018}
can be matched with an actual line bundle, not merely a numerical
class.  Taking the ample bundle $A$ of \eqref{geo:A-definition}, one
has, in additive notation,
\[
  A-\frac{A\cdot C}{C^2}C=A+C=2L_0.
\]
Thus the displayed construction in that paper includes
$L_0^{\otimes2}$.  The bundle $L_0$ itself is a convenient
distinguished choice on the same numerical ray.
\end{remark}

\subsection{The trace argument and failure of semiampleness}

\begin{lemma}\label{geo:no-compact-curves}
The affine bundle $T=X\setminus C$ contains no compact complex
analytic curve.
\end{lemma}

\begin{proof}
Suppose there were such a curve.  Passing to the reduced structure
and then to an irreducible component, choose a compact irreducible
curve $D\subset T$.  Let $\nu:\widetilde D\to D$ be its
normalization.  The resulting map
\[
  p:=q\circ\nu:\widetilde D\longrightarrow B
\]
cannot be constant.  Otherwise $D$ would lie in one affine fiber
$q^{-1}(b)\cap T\cong\C$; its affine coordinate would pull back to
a holomorphic function on the compact connected Riemann surface
$\widetilde D$, hence would be constant, forcing $D$ to be a point.

A nonconstant holomorphic map between compact connected Riemann
surfaces is surjective and finite.  Let its degree be $n\geq1$.
Away from the finitely many branch values of $p$, define on $U_i$
\begin{equation}\label{geo:trace-section}
  s_i(b)=\frac{1}{n}
          \sum_{x\in p^{-1}(b)}t_i(\nu(x)).
\end{equation}
There are $n$ points in this sum over an unbranched fiber.  Locally
they vary holomorphically.  Permuting the sheets does not change
their sum, so $s_i$ is a single-valued holomorphic function away
from the branch values.

We spell out why branching causes no problem.  Around a branch
value $b_1\in U_i$, choose a small closed disk contained in $U_i$.
Its inverse image in $\widetilde D$ is compact.  The holomorphic
function $t_i\circ\nu$ is bounded there, so the averages in
\eqref{geo:trace-section} remain bounded as $b$ approaches $b_1$.
The removable-singularity theorem extends $s_i$ holomorphically
across $b_1$.  This proves that $s_i$ is holomorphic on all of $U_i$.

On an unbranched fiber over $U_i\cap U_j$, summing
\eqref{geo:affine-transition} over its $n$ points gives
\[
  s_i=g_{ij}s_j+a_{ij}.
\]
The identity extends across branch values by holomorphicity.  Thus
the $s_i$ define a holomorphic section of $T\to B$, contradicting
Lemma~\ref{geo:no-section}.
\end{proof}

\begin{proposition}\label{geo:not-semiample}
For every integer $m>0$, every global section of $L_0^{\otimes m}$
vanishes identically on $C$.  In particular, $L_0$ is not semiample.
Its stable base locus is exactly $C$.
\end{proposition}

\begin{proof}
Since $L_0|_C\cong\OO_C$, the restriction of a section
$s\in H^0(X,L_0^{\otimes m})$ is a constant in the fixed
trivialization.  If this constant were nonzero, $s$ would be
nowhere zero on $C$.  Its zero divisor would therefore be disjoint
from $C$.

That zero divisor is nonempty.  Indeed, every fiber meets $C$ in
one point, so the restriction of $s$ to any fiber would be a
nonzero section of $\OO_{\P^1}(m)$.  Such a section has zeros
because $m>0$.  A nonzero holomorphic section of a line bundle
on a smooth surface has a divisor of zeros; here this is a
nonempty compact curve in $X\setminus C$.  This contradicts
Lemma~\ref{geo:no-compact-curves}.  Hence $s|_C=0$ for every $s$.

No positive tensor power of $L_0$ is globally generated, which is
precisely the failure of semiampleness.  On the other hand,
$A^{\otimes k}$ is globally generated for all sufficiently large
$k$, since $A$ is ample.  For each $x\notin C$, choose a section of
$A^{\otimes k}$ nonzero at $x$ and multiply it by $s_C^k$.
The resulting section of $L_0^{\otimes2k}$ is nonzero at $x$.
Thus no point outside $C$ belongs to the stable base locus, while
the first part shows that every point of $C$ does.
\end{proof}

\subsection{The canonical bundle and the contraction}

\begin{proposition}\label{geo:canonical}
The canonical bundle of $X$ is
\begin{equation}\label{geo:canonical-formula}
  K_X\cong\OO_X(-2C)\otimes q^*(K_B\otimes F^{-1}).
\end{equation}
In particular,
\[
  K_X\cdot C=2g-2+d>0,
\]
and $H^0(X,K_X^{\otimes m})=0$ for every $m>0$.  The smooth
surface $X$ therefore has Kodaira dimension $-\infty$.
\end{proposition}

\begin{proof}
Let $S=\OO_{\P_B(E)}(-1)$ be the tautological subbundle of $q^*E$.
The quotient $Q=q^*E/S$ has rank one.  Since $\det E=F$, taking
determinants gives $Q\cong q^*F\otimes S^{-1}$.  The relative
 tangent bundle of a projective bundle parametrizes first-order
deformations of a line in its fixed ambient fiber.
At a line $\ell\subset E_b$, a first-order deformation is the graph
of a linear map $\ell\to E_b/\ell$. These identifications vary
holomorphically, giving
\[
  T_{X/B}\cong\Hom(S,Q)
            \cong\OO_X(2C)\otimes q^*F.
\]
Dualizing and tensoring with $q^*K_B$ proves
\eqref{geo:canonical-formula}.  Intersecting with $C$ gives
$2d+(2g-2-d)=2g-2+d$, also in agreement with adjunction.

On every fiber $f\cong\P^1$, one has
$K_X^{\otimes m}|_f\cong\OO_{\P^1}(-2m)$, which has no nonzero
holomorphic sections.  Every global pluricanonical section
therefore vanishes on every fiber, hence on all of $X$.
\end{proof}

The next existence statement uses Grauert's analytic contraction
theorem as a classical background result, rather than an algebraic
base-point-free theorem.  In the local form needed here, that
theorem says that a connected compact curve in a smooth complex
surface whose intersection matrix is negative definite has a
neighborhood admitting a proper holomorphic contraction onto a
normal complex surface.  The curve is mapped to one point and the
map is biholomorphic off the curve
\cite{Grauert1962,BHPV2004}.

\begin{proposition}\label{geo:contraction}
There is a proper holomorphic map
\[
  \pi:X\longrightarrow Y
\]
onto a normal compact complex surface such that
$\pi(C)=\{y\}$ and $\pi$ is biholomorphic from $X\setminus C$
onto $Y\setminus\{y\}$.  The point $y$ is singular, and it is the
only singular point of $Y$.  The surface $Y$ is not projective.
\end{proposition}

\begin{proof}
The intersection matrix of the single curve $C$ is the negative
one-by-one matrix $(-d)$, so Grauert's theorem applies.  Let
$\pi_U:U\to V$ be its local contraction, with $V$ normal and
$\pi_U(C)=\{y\}$.  Glue $V$ to $X\setminus C$ using the
biholomorphism $U\setminus C\cong V\setminus\{y\}$.
Topologically the resulting space is the quotient obtained by
collapsing the compact closed subset $C$ to a point.  The proper
local contraction induces the quotient topology near that point.
This quotient is Hausdorff and compact, and its two indicated
charts supply a complex-space structure.  Normality holds near
$y$ by the local theorem and elsewhere because $X\setminus C$
is smooth.  The induced map $\pi$ is holomorphic and proper;
properness also follows directly from compactness of $X$ and
the Hausdorff property of $Y$.

We verify that $y$ cannot be smooth.  If it were, choose local
coordinates $z_1,z_2$ near $y$.  The pulled-back two-form
\[
  \eta=\pi^*(dz_1\wedge dz_2)
\]
would be a holomorphic section of $K_X$ over the inverse image
of that coordinate neighborhood.  It is nowhere zero off $C$,
since $\pi$ is biholomorphic there.  It vanishes on $C$, because
$\pi|_C$ is constant and therefore $d\pi$ annihilates its tangent
direction.  Its divisor is consequently $aC$ for an integer
$a\geq1$.  Restricting the corresponding line-bundle identity to
$C$ gives
\[
  K_X\cdot C=aC^2=-ad<0.
\]
This contradicts Proposition~\ref{geo:canonical}, or directly
the adjunction identity $K_X\cdot C=2g-2+d$.  Thus $y$ is singular.
All other points are smooth by the stated biholomorphism.

Finally, suppose $Y$ were projective.  In a projective embedding,
choose a hyperplane that avoids $y$ and does not contain $Y$.
Its intersection with the projective surface $Y$ is a nonempty
compact curve, disjoint from $y$.  The inverse image of that
curve under $\pi$ is a compact curve in $X\setminus C=T$,
contradicting Lemma~\ref{geo:no-compact-curves}.  Therefore $Y$
is not projective.
\end{proof}

\begin{remark}\label{geo:scope-singular}
The surface $Y$ is bimeromorphic to the projective surface $X$,
and hence is Moishezon.  The preceding argument establishes its
nonprojectivity directly.  We do not infer non-K\"ahlerness from a
general assertion about singular Moishezon spaces.  A separate
argument using the metric obstruction will establish that property
later.  Likewise, the Kodaira-dimension statement above concerns
the smooth ruled surface $X$; no assertion about reflexive powers
of the canonical sheaf on the singular space $Y$ is needed here.
\end{remark}

% End geometry.tex
% Begin obstruction.tex
\section{The obstruction to continuously differentiable potentials}
\label{obs:section}

We now prove that the regularity of the metric constructed in
Section~\ref{sec:metrics} cannot be
improved to $C^1$.  The argument is local near the infinity section, although
the compactness of that section is essential.  Its analytic ingredient is a
first-derivative property of plurisubharmonic functions.  In particular, we
will not differentiate a $C^1$ weight twice or assume that its curvature has
continuous coefficients.

\subsection{Regularity and the first normal derivative}

For a positive Hermitian metric $h$ on a holomorphic line bundle, write
\[
   |e|_h^2=e^{-\varphi}
\]
in a local holomorphic frame $e$.  We say that $h$ is of class $C^1$ if its
local weights $\varphi$ are of class $C^1$ as real-valued functions.  If
$e'=G e$, where $G$ is a nowhere-vanishing holomorphic function, then
\[
   \varphi'=\varphi-\log|G|^2.
\]
The change of weight is smooth, so this definition is independent of the
holomorphic frame.  The regularity of the metric coefficients is also
unchanged under a smooth change of frame.  Plurisubharmonicity of the
weights, however, is asserted in holomorphic frames.  A $C^1$ semipositive
metric means that these weights are plurisubharmonic, or equivalently that
their distributional $\ddc$ is positive.

We use the Wirtinger derivatives
\[
  \partial_z=\frac12(\partial_x-\sqrt{-1}\partial_y),
  \qquad
  \partial_{\bar z}=\frac12(\partial_x+\sqrt{-1}\partial_y),
  \qquad \Delta=4\partial_z\partial_{\bar z}.
\]

\begin{lemma}[Holomorphicity of the first normal derivative]
\label{obs:first-jet}
Let $D\subset\C$ be a disc, let $\Omega\subset\C^2$ be an open
neighborhood of $D\times\{0\}$, and let
$\varphi\in C^1(\Omega)\cap\PSH(\Omega)$ be real-valued.  Suppose that
$\varphi(z,0)$ is constant on $D$.  Then
\[
                  b(z):=\partial_w\varphi(z,0)
\]
is holomorphic on $D$.
\end{lemma}

\begin{proof}
Subtracting a constant, we may assume that $\varphi(z,0)=0$.  The function
$b$ is initially only known to be continuous.  We first record carefully
the uniformity of the first-order expansion.  Let $K\subset D$ be compact
and choose $\rho>0$ such that a neighborhood of
$K\times\{|w|\leq\rho\}$ is contained in $\Omega$.  The fundamental theorem
of calculus, applied to the real parameter $s$, gives
\begin{align*}
 \varphi(z,w)
 &=\int_0^1\frac{d}{ds}\varphi(z,sw)\,ds\\
 &=2\operatorname{Re}\left(
       w\int_0^1\partial_w\varphi(z,sw)\,ds\right).
\end{align*}
Here we used that $\varphi$ is real-valued.  Consequently
\begin{equation}
\label{obs:uniform-expansion}
 \varphi(z,w)=2\operatorname{Re}\bigl(b(z)w\bigr)+R(z,w),
 \qquad
 \sup_{z\in K}|R(z,w)|\leq 2|w|\omega_K(|w|),
\end{equation}
where one may take
\[
 \omega_K(\varepsilon)=
 \sup_{\substack{z\in K\\ |v|\leq\varepsilon}}
 \bigl|\partial_w\varphi(z,v)-\partial_w\varphi(z,0)\bigr|.
\]
Continuity of the first derivative on a compact neighborhood implies
$\omega_K(\varepsilon)\longrightarrow0$ as $\varepsilon\downarrow0$.
Thus the remainder is $o(|w|)$ locally uniformly in the base variable.

Let $f$ be holomorphic on $D$, and work first on a relatively compact
subdisc $D'\Subset D$.  For all sufficiently small real $t>0$, the two maps
\[
      z\longmapsto (z,t f(z)),\qquad
      z\longmapsto (z,-t f(z))
\]
take a neighborhood of $\overline{D'}$ into $\Omega$ and are holomorphic.
The functions
\[
       u_t^+(z)=\frac{\varphi(z,t f(z))}{t},
       \qquad
       u_t^-(z)=\frac{\varphi(z,-t f(z))}{t}
\]
are therefore subharmonic on that neighborhood.  The denominator is
positive in both expressions; no multiplication of a subharmonic function
by a negative constant is being used.  Estimate
\eqref{obs:uniform-expansion} proves the locally uniform limits
\begin{equation}
\label{obs:graph-limits}
       u_t^+\longrightarrow 2\operatorname{Re}(bf),
       \qquad
       u_t^-\longrightarrow-2\operatorname{Re}(bf).
\end{equation}
For example, on a compact subset where $|f|\leq A$, the error in either
limit is bounded by $2A\omega_K(tA)$.

For completeness, these limits are subharmonic directly from the
submean inequality.  If a closed disc lies in $D'$, the inequality for
$u_t^\pm$ compares its value at the center with its circular mean.
Uniform convergence on that disc and its boundary permits passage to
the limit.  The resulting functions are continuous and satisfy the same
inequality on every such disc, hence are subharmonic.  It follows from
\eqref{obs:graph-limits} that both
$\operatorname{Re}(bf)$ and its negative are subharmonic.

We also recall why this implies ordinary harmonicity, with no regularity
assumption beyond continuity.  Applying the submean inequality to a
continuous function and to its negative gives equality with every
sufficiently small circular mean.  Convolution with a smooth radial
mollifier is a weighted average of these circular means and therefore
equals the function wherever the convolution is defined.  The function
is consequently smooth there.  Its circular-mean identity, expanded to
second order, then gives $\Delta h=0$.  Thus
$\operatorname{Re}(bf)$ is harmonic on $D'$, and hence on $D$ because
$D'$ was arbitrary.

Take successively $f=1$, $f=\sqrt{-1}$, $f=z$, and
$f=\sqrt{-1}z$.  The first two choices show that the real and imaginary
parts of $b$ are harmonic; the last two show the same for $zb$.  In
particular $b$ is smooth and
\[
                    \Delta b=0,\qquad \Delta(zb)=0.
\]
The product rule for Wirtinger derivatives gives the exact identity
\[
   \Delta(zb)
     =4\partial_z\bigl(z\partial_{\bar z}b\bigr)
     =4\partial_{\bar z}b+z\Delta b.
\]
Hence $\partial_{\bar z}b=0$, which proves that $b$ is holomorphic.
\end{proof}

\begin{remark}
\label{obs:c2-comparison}
There is a shorter proof when $\varphi$ is $C^2$.  Its complex Hessian at
$(z,0)$ is a positive semidefinite Hermitian matrix, and its
$z\bar z$ entry vanishes because $\varphi(z,0)=0$.  Positivity of a
two-by-two Hermitian matrix with one zero diagonal entry forces the
corresponding mixed entries to vanish.  Thus
$\varphi_{w\bar z}(z,0)=0$, and
$\partial_{\bar z}b=0$.  Lemma~\ref{obs:first-jet} replaces this Hessian
argument when only one continuous derivative is available.
\end{remark}

\subsection{The transition calculation on the infinity section}

We spell out the conventions needed for the application.  Let
$q:X\longrightarrow B$ be the ruled surface obtained by compactifying the
affine line bundle associated with a nonzero class
$\xi\in H^1(B,F)$, and let $C\subset X$ be its infinity section.  Choose
local holomorphic frames $m_i$ of $F$ and affine fiber coordinates $t_i$
over open sets $U_i\subset B$, with transition formulas
\begin{equation}
\label{obs:affine-transitions}
       m_j=g_{ij}m_i,
       \qquad t_i=g_{ij}t_j+a_{ij}.
\end{equation}
The functions $g_{ij}$ are nowhere vanishing and holomorphic, the
functions $a_{ij}$ are holomorphic, and the cocycle
$\{a_{ij}m_i\}$ represents $\xi$.  Base points are suppressed in these
formulas; a function on $U_j$ is understood to be restricted to the
overlap and expressed in the chosen base coordinate there.

A holomorphic section of the affine line bundle would be given by
holomorphic functions $s_i$ satisfying
\begin{equation}
\label{obs:section-equation}
                    s_i=g_{ij}s_j+a_{ij}.
\end{equation}
Multiplying this identity by $m_i$ gives
$s_i m_i-s_jm_j=a_{ij}m_i$.  Thus such functions would express the
cocycle as a coboundary.  Since $\xi\ne0$, no such section exists.

Put
\[
                  L_0=\OO_X(C)\otimes q^*F.
\]
Near $C$, use $w_i=1/t_i$ as the transverse coordinate, so that
$C=\{w_i=0\}$.  Formula~\eqref{obs:affine-transitions} becomes
\begin{equation}
\label{obs:normal-transitions}
   w_i=\frac{w_j}{g_{ij}+a_{ij}w_j},
   \qquad
   w_j=\frac{g_{ij}w_i}{1-a_{ij}w_i}.
\end{equation}
All expressions here are considered on sufficiently small overlapping
neighborhoods of $C$, where their denominators do not vanish.  If
$\sigma_C$ is the canonical section of $\OO_X(C)$, choose its natural
local frames $v_i$ by
\[
                        \sigma_C=w_i v_i.
\]
The identity $w_i v_i=w_jv_j$ and
\eqref{obs:normal-transitions} imply
\[
                  v_j=\frac{1-a_{ij}w_i}{g_{ij}}v_i.
\]
Consequently the holomorphic frames
$e_i=q^*m_i\otimes v_i$ of $L_0$ satisfy the exact formula
\begin{equation}
\label{obs:line-transitions}
                         e_j=(1-a_{ij}w_i)e_i.
\end{equation}
In particular, their restrictions agree on $C$ and define a global
nowhere-vanishing holomorphic section of $L_0|_C$.

\begin{theorem}[Nonexistence of a $C^1$ semipositive metric]
\label{obs:c1-obstruction}
Assume that the affine extension class $\xi\in H^1(B,F)$ is nonzero.
Then $L_0$ admits no $C^1$ Hermitian metric with semipositive curvature
current on any open neighborhood of $C$ in $X$.
\end{theorem}

\begin{proof}
Suppose that such a metric $h$ exists on a neighborhood $N$ of $C$.
Shrink the local coordinate neighborhoods used above so that they lie
in $N$, and write
\[
                        |e_i|_h^2=e^{-\varphi_i}.
\]
The functions $\varphi_i$ are $C^1$ and plurisubharmonic.  The exact
transformation law~\eqref{obs:line-transitions} gives
\begin{equation}
\label{obs:weight-transitions}
 \varphi_j(z,w_j)
    =\varphi_i(z,w_i)-\log|1-a_{ij}(z)w_i|^2,
 \qquad
 w_j=\frac{g_{ij}(z)w_i}{1-a_{ij}(z)w_i}.
\end{equation}
Here the same notation $z$ denotes the base point on the overlap; any
change between base coordinate charts depends only on that base point.

At $w_i=w_j=0$, the logarithmic term vanishes.  Thus the restrictions
$\varphi_i|_C$ glue to a global continuous real-valued function on $C$.
The restriction of a finite plurisubharmonic function to a holomorphic
curve is subharmonic.  Since $C$ is compact and connected, the maximum
principle implies that this global function is constant.  Subtracting
the same constant from all weights, we arrange that
\begin{equation}
\label{obs:zero-restriction}
                         \varphi_i(z,0)=0.
\end{equation}
This normalization does not change curvature or regularity.

Set
\[
                       b_i(z)=\partial_{w_i}\varphi_i(z,0).
\]
Lemma~\ref{obs:first-jet}, applied on coordinate discs in $U_i$, proves
that every $b_i$ is holomorphic.  We may differentiate
\eqref{obs:weight-transitions} once in $w_i$, because the weights are
$C^1$.  At $w_i=0$, the derivative of $w_j$ with respect to $w_i$ is
$g_{ij}$, and its antiholomorphic counterpart contributes nothing to
the Wirtinger derivative.  Moreover,
\[
 \left.\partial_{w_i}\log|1-a_{ij}w_i|^2\right|_{w_i=0}
                  =-a_{ij}.
\]
We therefore obtain
\begin{equation}
\label{obs:first-jet-transition}
                    g_{ij}b_j=b_i+a_{ij},
                    \qquad b_i=g_{ij}b_j-a_{ij}.
\end{equation}
Equivalently, one can substitute the locally uniform expansions from
\eqref{obs:uniform-expansion} into
\eqref{obs:weight-transitions} and compare the real linear terms in
$w_i$; this gives the same identity.

Now put $s_i=-b_i$.  The functions $s_i$ are holomorphic, and
\eqref{obs:first-jet-transition} becomes precisely
\[
                         s_i=g_{ij}s_j+a_{ij}.
\]
They define the forbidden affine section from
\eqref{obs:section-equation}.  Explicitly,
\[
                  b_jm_j-b_im_i=a_{ij}m_i
\]
expresses the cocycle representing $\xi$ as a coboundary.  This
contradicts $\xi\ne0$ and proves the theorem.
\end{proof}

\begin{corollary}
\label{obs:powers}
For every integer $m\geq1$, the line bundle $L_0^{\otimes m}$ admits
no $C^1$ semipositive Hermitian metric on a neighborhood of $C$.
Equivalently, relative to any smooth reference metric on
$L_0^{\otimes m}$, no semipositive metric potential can be $C^1$ on a
neighborhood of $C$.
\end{corollary}

\begin{proof}
A positive metric on $L_0^{\otimes m}$ has a unique positive $m$th root
as a metric on the already given line bundle $L_0$.  In the frames
$e_i^{\otimes m}$ and $e_i$, this simply divides every local weight by
$m$.  It preserves $C^1$ regularity and plurisubharmonicity, contradicting
Theorem~\ref{obs:c1-obstruction}.  Subtracting the smooth local weight
of a reference metric does not change whether a potential is $C^1$.
\end{proof}

\subsection{Relation with the first obstruction class}

\begin{remark}
\label{obs:obstruction-class}
The calculation also identifies the first obstruction in the
transition-function approach.  With the convention that this
obstruction is represented by the coefficient of $w_i$ in $e_j/e_i$,
formula~\eqref{obs:line-transitions} gives
\[
     u_1(C,X,L_0)
        =\bigl[\{-a_{ij}\,dw_i\}\bigr]
           \in H^1(C,N_{C/X}^{-1}).
\]
Along $C$, formula~\eqref{obs:normal-transitions} yields
$dw_j=g_{ij}dw_i$.  Hence the identification
$N_{C/X}^{-1}\simeq F$ sending $dw_i$ to $m_i$ identifies this class
with $-\xi$.  Reversing the convention for the transition cocycle
changes the sign but not its nonvanishing.  For a tensor power, the
expansion
\[
       \frac{e_j^{\otimes m}}{e_i^{\otimes m}}
           =(1-a_{ij}w_i)^m
           =1-ma_{ij}w_i+O(w_i^2)
\]
identifies the corresponding class with $-m\xi$.

Koike's line-bundle obstruction theory~\cite{Koike2021} extends the
transition-function viewpoint originating in Ueda's neighborhood
theory~\cite{Ueda1983}.  Zhang computes the nonzero first obstruction
for the original Grauert construction and applies the smooth
semipositivity obstruction~\cite{Zhang2025}.  The proof above is given
independently: its additional analytic ingredient is
Lemma~\ref{obs:first-jet}, and it establishes nonexistence for $C^1$
weights without invoking a theorem that assumes a smooth metric.
\end{remark}

\begin{remark}[What the regularity conclusion means]
\label{obs:regularity-meaning}
The obstruction is to continuous first derivatives, not to bounded or
Lipschitz local weights.  A Lipschitz function need not possess the
uniform real-linear first-order expansion
\eqref{obs:uniform-expansion} along $C$.  In particular, the transverse
model $|w|$ is Lipschitz but is not differentiable at $w=0$.  This is
exactly the type of regularity permitted by the norm-of-section
potential in the construction.  Thus Lipschitz existence and
Theorem~\ref{obs:c1-obstruction} are compatible.  Together they give
the precise endpoint statement used here: Lipschitz semipositive
potentials exist, whereas $C^1$ semipositive potentials do not.  No
classification of other, weaker differentiability notions is asserted.
\end{remark}

% End obstruction.tex
% Begin metrics.tex
\section{Constructing a Lipschitz semipositive metric}
\label{sec:metrics}

The construction in this section applies more generally than to a ruled
surface. Its local input is a negative smooth curve and a line bundle
whose restriction to that curve is flat. The global input is an exact
ample twist by the curve. We prove the local estimates and the gluing
separately, using the same smooth reference metric throughout.

\subsection{Reference metrics near a negative curve}

Let $S$ be a smooth complex surface and let $C\subset S$ be a smooth,
compact, connected complex curve with $C^2<0$. Let $L$ be a holomorphic
line bundle on a neighborhood of $C$. Assume that $L|_C$ has a smooth
Hermitian metric of zero curvature. In our application $L|_C$ is
holomorphically trivial, so the constant metric in a global nonvanishing
frame has this property.

We explain two choices of smooth metrics. First choose a smooth metric
$h_0$ on $L$ whose restriction is the specified flat metric. To do this,
start with any smooth metric $h'_0$. The logarithm of the ratio of the
desired metric to $h'_0|_C$ is a smooth real function on $C$. Extend that
function smoothly to a neighborhood of $C$, and multiply $h'_0$ by its
exponential. Smooth extension is available from a smooth tubular
neighborhood, or from smooth coordinate extensions and a partition of
unity. No holomorphic tubular neighborhood is used. With
\[
                  \theta=\cOne(L,h_0),
\]
we have $\theta|_{TC}=0$.

Second, choose a smooth Hermitian metric $k$ on $\OO_S(C)$ such that
\begin{equation}\label{met:negative-metric}
       \eta:=-\cOne(\OO_S(C),k)
       \quad\hbox{is strictly positive on }TC.
\end{equation}
Here is why such a choice is possible. The restriction
$\OO_S(C)|_C=N_{C/S}$ has degree $C^2<0$. On a compact Riemann surface,
a line bundle of negative degree has a smooth metric with strictly
negative curvature. In detail, fix a positive area form $\omega_C$ of
integral one and any metric with curvature $\gamma$. The form
$C^2\omega_C-\gamma$ has integral zero. The Poisson equation on the
compact curve gives a smooth real function $v$ with
$\ddc v=C^2\omega_C-\gamma$. Multiplying the metric by $e^{-v}$
gives curvature $C^2\omega_C$. The Poisson solvability used here is the
usual compact-surface statement: the Laplacian has constants as kernel
and has as range the smooth functions of integral zero
\cite{DemaillyCADG}. Extend this metric from $C$ to $\OO_S(C)$ by the
smooth ratio construction just described. Its negative curvature has
the required positive restriction.

Let $s_C$ be the canonical section of $\OO_S(C)$ and set
\begin{equation}\label{met:r-s}
                    r=|s_C|_k^2,\qquad s=r^{1/2}.
\end{equation}
Thus $r$ is a smooth nonnegative function, with zero set $C$. In a
holomorphic coordinate chart $(z,w)$ with $C=\{w=0\}$, choose a frame
$v$ of $\OO_S(C)$ such that $s_C=w v$ and put
$|v|_k^2=e^{\lambda}$. Then
\begin{equation}\label{met:r-coordinates}
       r=|w|^2e^{\lambda},\qquad
       s=|w|e^{\lambda/2},\qquad \eta=\ddc\lambda.
\end{equation}
In particular, $s$ is locally Lipschitz, including on $C$, and is smooth
off $C$. Indeed, $w\mapsto|w|$ is Lipschitz, and $e^{\lambda/2}$ is
smooth with bounded first derivatives on every relatively compact
coordinate chart; their product is Lipschitz there.

On $S\setminus C$ one has
\begin{equation}\label{met:log-r}
                       \ddc\log r=\eta.
\end{equation}
This is an identity on the complement only. As currents across $C$,
the correct identity is
\begin{equation}\label{met:PL-r}
                       \ddc\log r=[C]+\eta.
\end{equation}
It follows from $\ddc\log|w|^2=[w=0]$ with our normalization. This
distinction will be essential for the global auxiliary metric.

\subsection{A frame adapted to the level sets of the norm}

We will estimate Hermitian coefficient matrices as in
Section~\ref{sec:preliminaries}: a real $(1,1)$-form is written as
$(\sqrt{-1}/2\pi)\sum T_{\alpha\bar\beta}
 dz_\alpha\wedge d\bar z_\beta$, and positivity is tested on
$(T_{\alpha\bar\beta})$. All coefficients below use this convention.

Shrink a relatively compact coordinate chart until
$|w\lambda_w|\le1/2$. Define two smooth $(1,0)$ vector fields
\begin{equation}\label{met:adapted-frame}
 H=\partial_z-\frac{w\lambda_z}{1+w\lambda_w}\partial_w,
 \qquad V=\partial_w.
\end{equation}
They form a frame because the coefficient of $\partial_z$ in $H$ is
one. Its change-of-frame matrix is triangular with determinant one.
It extends smoothly to $w=0$, where $H=\partial_z$.
On $w\ne0$, direct substitution gives
\begin{align}
 \partial\log r(H)
 &=\lambda_z-
   \frac{w\lambda_z}{1+w\lambda_w}
                    \left(\frac1w+\lambda_w\right)=0,
                    \label{met:H-log}\\
 \partial\log r(V)&=\frac1w+\lambda_w.
                    \label{met:V-log}
\end{align}
The vector fields need not be holomorphic. We use them only to
\emph{evaluate} tensors computed in the holomorphic coordinates.
In particular, none of the formulas below replaces a Hessian coefficient
by a repeated derivative along the nonholomorphic field $H$.

For $\delta>0$, differentiating $r^\delta=e^{\delta\log r}$ on the
complement of $C$ gives the exact formula
\begin{equation}\label{met:power-hessian}
 \ddc r^\delta
   =\delta r^\delta\eta+
     \delta^2r^\delta\frac{\sqrt{-1}}{2\pi}
                     \partial\log r\wedge\bar\partial\log r.
\end{equation}
In the frame $(H,V)$ this reads
\begin{align}
 (\ddc r^\delta)_{H\bar H}
      &=\delta r^\delta\eta_{H\bar H},\label{met:HH-exact}\\
 (\ddc r^\delta)_{H\bar V}
      &=\delta r^\delta\eta_{H\bar V},\label{met:HV-exact}\\
 (\ddc r^\delta)_{V\bar V}
      &=\delta r^\delta\eta_{V\bar V}
       +\delta^2r^\delta\left|\frac1w+\lambda_w\right|^2.
          \label{met:VV-exact}
\end{align}
The cancellation in \eqref{met:HV-exact} is why we use the adapted
frame. Estimating entries independently in the original frame would
hide the useful rank-one structure of the second term in
\eqref{met:power-hessian}.

\begin{lemma}[Uniform estimates for the square-root cusp]
\label{met:cusp-estimates}
There are a neighborhood $U$ of $C$, finitely many charts covering
$U$, and constants $a,b,K>0$ with the following properties on
$U\setminus C$. In each chart, using the frame
\eqref{met:adapted-frame},
\begin{align}
 (\ddc s)_{H\bar H}&\ge a s,&
 (\ddc s)_{V\bar V}&\ge \frac b{s},&
 |(\ddc s)_{H\bar V}|&\le K s,\label{met:cusp-bounds}\\
 \theta_{H\bar H}&\ge-Ks,&
 \theta_{V\bar V}&\ge-K,&
 |\theta_{H\bar V}|&\le K.\label{met:theta-bounds}
\end{align}
All constants are independent of the point and of the chart in this
finite covering.
\end{lemma}

\begin{proof}
Choose finitely many relatively compact charts around $C$ and smaller
charts whose intersections with $C$ still cover it. Work on these
smaller charts, so all smooth functions and their needed derivatives
have bounds supplied by the larger charts. After shrinking in the
normal direction, $|w\lambda_w|\le1/2$ holds on every chart. Since
$H|_C=\partial_z$ and $\eta|_{TC}>0$, continuity and the finite cover
give a number $\mu>0$ such that
$\eta_{H\bar H}\ge\mu$ on a common smaller neighborhood. There is
also a common $A\ge1$ such that
\[
 |\eta_{H\bar V}|\le A,\qquad |\eta_{V\bar V}|\le A,
 \qquad A^{-1}\le e^\lambda\le A.
\]
Put $\delta=1/2$ in \eqref{met:HH-exact}--\eqref{met:VV-exact}.
The $H\bar H$ and $H\bar V$ formulas immediately give
\[
 (\ddc s)_{H\bar H}\ge\frac\mu2 s,
 \qquad |(\ddc s)_{H\bar V}|\le\frac A2s.
\]
The bound $|w\lambda_w|\le1/2$ implies
\[
 \left|\frac1w+\lambda_w\right|
   =\frac{|1+w\lambda_w|}{|w|}\ge\frac{1}{2|w|}.
\]
Using $s^2=|w|^2e^\lambda$, we obtain
\begin{align*}
 (\ddc s)_{V\bar V}
 &\ge-\frac A2s+\frac{s}{16|w|^2}
  =-\frac A2s+\frac{e^\lambda}{16s}\\
 &\ge-\frac A2s+\frac{1}{16A s}.
\end{align*}
Shrink the neighborhood so that $s^2\le1/(16A^2)$.
Then $(A/2)s\le1/(32A s)$, and the last expression is at least
$1/(32A s)$. We may take $a=\mu/2$ and $b=1/(32A)$.

For the reference form, $\theta_{H\bar H}$ is a smooth real function
on the entire chart, including $w=0$. Its value there is zero because
$H|_C=\partial_z$ and $\theta|_{TC}=0$. The real mean-value estimate
in the two normal variables therefore gives
$|\theta_{H\bar H}(z,w)|\le A_1|w|$ for a uniform constant $A_1$.
For example, integrate its first normal derivative along the segment
$t\mapsto(z,tw)$, $0\le t\le1$, in a product subchart. Since
$|w|\le A^{1/2}s$, this is at most $A_1A^{1/2}s$. The other two
coefficients of $\theta$ in the smooth frame are uniformly bounded.
Enlarge $K$ to dominate those bounds, $A_1A^{1/2}$, and $A/2$.
The displayed inequalities follow on the common neighborhood.
\end{proof}

\subsection{Fixing the coefficient before shrinking the neighborhood}

\begin{theorem}[Local Lipschitz existence]
\label{met:lipschitz-local}
Let $C$ and $L$ satisfy the assumptions at the beginning of this
section. There exist a neighborhood $U$ of $C$ and $M>0$ such that
\begin{equation}\label{met:local-current}
                    T_M=\theta+M\ddc s
\end{equation}
is a positive current on $U$ and a smooth strictly positive
$(1,1)$-form on $U\setminus C$. The metric
$h_{\mathrm{loc}}=h_0e^{-Ms}$ is positive definite with locally
Lipschitz weights on $U$.
\end{theorem}

\begin{proof}
Choose $a,b,K$ as in Lemma~\ref{met:cusp-estimates}. First choose one
constant $M>0$ satisfying
\begin{equation}\label{met:choose-M}
       Ma-K>0,\qquad
       (Ma-K)\frac{Mb}{2}>(K+1)^2.
\end{equation}
Such a choice exists: the left side of the second inequality is a
quadratic polynomial in $M$ with positive leading coefficient
$ab/2$. This choice is made before any further shrinking.

With this $M$ fixed, shrink $U$ until
\begin{equation}\label{met:shrink-after-M}
              Ks\le\frac{Mb}{2},\qquad MKs\le1
\end{equation}
throughout $U$. This is possible because $s$ is continuous and zero
on the compact set $C$. The estimates in the preceding lemma now give
\begin{align*}
 (T_M)_{H\bar H}&\ge(Ma-K)s,\\
 (T_M)_{V\bar V}&\ge\frac{Mb}{s}-K
                         \ge\frac{Mb}{2s},\\
 |(T_M)_{H\bar V}|&\le K+MKs\le K+1.
\end{align*}
Both diagonal coefficients are strictly positive for $s>0$. Moreover,
\begin{align*}
 \det (T_M)_{(H,V)}
 &= (T_M)_{H\bar H}(T_M)_{V\bar V}
                         -|(T_M)_{H\bar V}|^2\\
 &\ge (Ma-K)\frac{Mb}{2}-(K+1)^2>0.
\end{align*}
The two-by-two Hermitian matrix is therefore positive definite.
Because $(H,V)$ is a frame, this is positivity of the form itself.
This proves strict positivity on $U\setminus C$.

Let $e$ be any local holomorphic frame of $L$ and write
$|e|_{h_0}^2=e^{-\varphi_0}$. The weight
$\varphi_0+Ms$ is smooth and strictly plurisubharmonic off $C$ and
is continuous across $C$. The removable lemma
Lemma~\ref{pre:removable} makes it plurisubharmonic on the whole chart.
Its distributional $\ddc$ is thus positive on that chart, including
on $C$. These local statements give positivity of $T_M$ as a current
on $U$. Finally, $\varphi_0$ is smooth and $s$ is locally Lipschitz
by \eqref{met:r-coordinates}, so the weight has the asserted regularity.
It is finite everywhere; its exponential defines a positive definite
Hermitian metric, not a metric degenerating to zero or infinity on $C$.
\end{proof}

\begin{remark}\label{met:exponent-remark}
The coefficient choice in \eqref{met:choose-M} is particularly useful
at the Lipschitz exponent $s=r^{1/2}$. At that exponent the product of
the leading tangential and normal lower bounds is of constant order,
so merely shrinking the neighborhood need not dominate a fixed mixed
coefficient. Increasing $M$ first supplies that domination. This is
also why an argument that only says that the normal coefficient tends
to infinity would not suffice: positivity requires the determinant
inequality, not just the two diagonal inequalities.
\end{remark}

\subsection{The global auxiliary metric and its exact singularity}

For global existence, suppose now that $S$ is compact and projective,
$L$ is defined on $S$, and there are integers $m,a_0>0$ such that
\begin{equation}\label{met:ample-twist}
                     A=L^{\otimes m}\otimes\OO_S(-a_0C)
                     \quad\hbox{is ample}.
\end{equation}
All line-bundle identities in what follows are holomorphic identities,
not only numerical equivalences. Choose $h_0$ and $k$ globally with
the restrictions used above. Global extensions of the prescribed
metrics can again be made by a smooth extension and a cut-off equal
to one near $C$.

An ample bundle has a smooth strictly positive metric $h_A$. Explicitly,
some $A^{\otimes N}$ is very ample and gives an embedding into
projective space. Pull back the Fubini--Study metric on $\OO(1)$
and take its $N$th root. Its curvature is
$\omega_A=N^{-1}\iota^*\omega_{\mathrm{FS}}>0$.

The singular divisor metric $h_{[C]}$ on $\OO_S(C)$ is characterized
by $|s_C|_{h_{[C]}}^2=1$ away from $C$. If $s_C=w v$, its local
weight in the frame $v$ is $\log|w|^2$, so its curvature current is
$[C]$. Relative to $k$ it satisfies
\[
                         h_{[C]}=k e^{-\log r}.
\]
Using the isomorphism $L^{\otimes m}=A\otimes\OO_S(a_0C)$,
form the singular metric
\[
                         h_*^m=h_A\otimes h_{[C]}^{a_0}.
\]
This defines a metric on $L$: the positive $m$th root of the squared
norm on $L^{\otimes m}$ is independent of the choice of frame.
Its curvature current is
\begin{equation}\label{met:auxiliary-current}
            \cOne(L,h_*)=\frac1m\omega_A+\frac{a_0}{m}[C].
\end{equation}

For the gluing we need a \emph{global function} relative to $h_0$.
Define a smooth real function $F_0$ on $S$ by the ratio of two smooth
metrics on the same line bundle:
\[
                       h_A\otimes k^{a_0}=h_0^m e^{-F_0}.
\]
The ratio is a globally defined positive smooth function, so its
negative logarithm $F_0$ is globally well-defined. Then
\begin{equation}\label{met:psi}
   h_*=h_0e^{-\psi},\qquad
   \psi=\frac{a_0}{m}\log r+\frac{F_0}{m},\qquad
   \theta+\ddc\psi=\frac1m\omega_A+\frac{a_0}{m}[C].
\end{equation}
In particular, $\psi$ is smooth on $S\setminus C$, tends to $-\infty$
uniformly on approach to $C$, and is strictly $\theta$-plurisubharmonic
off $C$. The coefficient of its logarithmic singularity is positive.

\subsection{Gluing by a regularized maximum}

\begin{theorem}[Global Lipschitz existence]
\label{met:lipschitz-global}
Let $S$ be a smooth projective surface, let $C\subset S$ be a smooth
connected curve with $C^2<0$, and let $L$ be a holomorphic line bundle.
Suppose $L|_C$ admits a smooth flat Hermitian metric and
\eqref{met:ample-twist} holds for some positive integers $m,a_0$.
Then $L$ admits a positive definite Hermitian metric with locally
Lipschitz weights and semipositive curvature current. The metric can
be chosen smooth with strictly positive curvature on $S\setminus C$.
Its curvature has bounded local potentials and minimal singularities
in $\cOne(L)$.
\end{theorem}

\begin{proof}
Use Theorem~\ref{met:lipschitz-local} to choose $U$ and $M$ so that
\[
                         u_{\mathrm{loc}}=Ms
\]
has the asserted positivity on $U$. Choose a small number $b>0$
such that $\{r\le2b\}\subset U$. This is possible because $S$ is
compact, $r$ is continuous with zero set exactly $C$, and the compact
set $S\setminus U$ has a positive minimum of $r$ whenever it is
nonempty. Decrease $b$ further if necessary. Construct $\psi$ as in
\eqref{met:psi}, and fix a regularization size $\varepsilon>0$.

On the compact shell $\{b/2\le r\le b\}$ the functions $\psi$ and
$u_{\mathrm{loc}}$ are smooth. Choose a real constant $N$ such that
\begin{equation}\label{met:outer-collar}
                 \psi+N>u_{\mathrm{loc}}+2\varepsilon
                     \quad\hbox{on }\{b/2\le r\le b\}.
\end{equation}
For instance, any $N>2\varepsilon+
\max_{b/2\le r\le b}(u_{\mathrm{loc}}-\psi)$ works; if the shell
is empty the requirement is vacuous. The boundedness of $F_0$ and
the positive coefficient of $\log r$ imply that $\psi+N$ tends
uniformly to $-\infty$ as $r\downarrow0$. Since
$u_{\mathrm{loc}}\ge0$, there exists $a$ with $0<a<b/4$ such that
\begin{equation}\label{met:inner-collar}
                 \psi+N<u_{\mathrm{loc}}-2\varepsilon
                         \quad\hbox{on }\{0<r\le2a\}.
\end{equation}
For complete uniformity, if $D_0=\max_S F_0/m+N$, it suffices to
choose $a<b/4$ so that
$(a_0/m)\log(2a)+D_0<-2\varepsilon$.

Let $\mathcal M_\varepsilon$ be the smooth regularized maximum from
Lemma~\ref{pre:regularized-max}. Define a real function on $S$ by
\begin{equation}\label{met:global-potential}
 u(x)=
 \begin{cases}
    u_{\mathrm{loc}}(x),& r(x)\le a,\\
    \mathcal M_\varepsilon
             \bigl(u_{\mathrm{loc}}(x),\psi(x)+N\bigr),
                                    &a<r(x)<b,\\
    \psi(x)+N,&r(x)\ge b.
 \end{cases}
\end{equation}
On a whole collar of the inner interface $r=a$, condition
\eqref{met:inner-collar} makes the middle expression equal to
$u_{\mathrm{loc}}$. On a whole collar of the outer interface $r=b$,
condition \eqref{met:outer-collar} makes it equal to $\psi+N$.
Consequently the three pieces agree in neighborhoods of their
interfaces. No regularity of the level hypersurfaces is needed for
this pasting assertion.

We check positivity in local holomorphic frames of $L$. Let
$\varphi_0$ be the weight of $h_0$. On the annulus, translation
invariance of the regularized maximum gives
\begin{align*}
 \varphi_0+
 \mathcal M_\varepsilon(u_{\mathrm{loc}},\psi+N)
 =\mathcal M_\varepsilon
        (\varphi_0+u_{\mathrm{loc}},\varphi_0+\psi+N).
\end{align*}
Both input functions are smooth and strictly plurisubharmonic there:
the first by the local theorem and the second by
\eqref{met:auxiliary-current} off $C$. By
Lemma~\ref{pre:regularized-max} and Remark~\ref{pre:strict-max}, their regularized maximum
is smooth and strictly plurisubharmonic on the annulus. On the inner
region positivity follows from the local theorem. On the outer
region it follows from \eqref{met:auxiliary-current}. The collars
show these assertions fit together, proving that
$\theta+\ddc u$ is positive globally and strictly positive away
from $C$.

Near $C$, the function $u$ equals $Ms$, hence is locally Lipschitz.
On the complement it is smooth, including at both interfaces by the
collar equalities. Thus $u$ is a globally defined, locally Lipschitz,
finite function on compact $S$, and is bounded. The metric
$h=h_0e^{-u}$ is positive definite with locally Lipschitz weights;
it is smooth off $C$ and has curvature $\theta+\ddc u$.
Finally, Lemma~\ref{pre:bounded-minimal} proves minimal singularities and
bounded local potentials. This last step is a comparison with all
$\theta$-plurisubharmonic potentials, not a consequence of nefness
alone.
\end{proof}

\begin{remark}\label{met:no-partition}
The gluing combines potentials relative to the same $\theta$ and uses
the holomorphic identity \eqref{met:ample-twist}. A partition of unity
applied directly to plurisubharmonic weights would introduce uncontrolled
derivative terms. Likewise, a merely numerical equality with an ample
class would not by itself justify the metric identity used to define
$F_0$. Both issues are avoided by the explicit construction above.
\end{remark}

% End metrics.tex
% Begin consequences.tex
\section{Applications to the original and modified examples}
\label{sec:consequences}

\begin{corollary}\label{con:application}
For the ruled surface and the line bundle $L_0$ of
Assumption~\ref{geo:data}, there exists a positive definite Hermitian metric
on $L_0$ with locally Lipschitz weights and semipositive curvature
current. The metric is smooth and has strictly positive curvature
on $X\setminus C$, and its curvature has bounded local potentials
and minimal singularities. No neighborhood of $C$ admits a $C^1$
semipositive Hermitian metric on $L_0$.
\end{corollary}

\begin{proof}
The surface $X$ is smooth and projective by
Proposition~\ref{geo:projective-bundle}. The curve $C$ is smooth and connected,
and $C^2=-d<0$, by Proposition~\ref{geo:normal-and-L}. That proposition also
provides the holomorphic trivialization $L_0|_C\cong\OO_C$; its
constant Hermitian metric is flat. Finally,
Proposition~\ref{geo:nef-big-ample} proves that the actual holomorphic line
bundle
\[
                A=L_0^{\otimes2}\otimes\OO_X(-C)
\]
is ample. These are all the hypotheses of
Theorem~\ref{met:lipschitz-global}, with $S=X$, $L=L_0$, $m=2$, and
$a_0=1$. That theorem proves the asserted existence and all its
regularity and minimality properties. The class $\xi$ is nonzero
by assumption, so Theorem~\ref{obs:c1-obstruction} gives nonexistence at
the $C^1$ level, even on a neighborhood of $C$.
\end{proof}

\begin{proof}[Completion of the proof of Theorem~\ref{main:regularity}]
Part (i) is Propositions~\ref{geo:projective-bundle} and \ref{geo:normal-and-L}.
The intersection, nefness, bigness, and null-locus assertions in
part (ii) are Proposition~\ref{geo:nef-big-ample}, and the statement about all
positive powers is Proposition~\ref{geo:not-semiample}.
Parts (iii) and (iv) are precisely Corollary~\ref{con:application}.
\end{proof}

\subsection{The degree-one choice and the canonical-bundle choice}

For an entirely specified degree-one instance, let $B$ have genus
two, choose $b_0\in B$, and put $F=\OO_B(b_0)$.
Lemma~\ref{geo:degree-one-choice} gives
$H^1(B,F)\cong\C$, so choose any nonzero element. If
$f_{b_0}=q^{-1}(b_0)$ denotes the actual fiber divisor, then
\begin{equation}\label{con:degree-one}
 L_0=\OO_X(C+f_{b_0}),\qquad C^2=-1,\qquad L_0^2=1.
\end{equation}
The line-bundle equality in this display uses $F=\OO_B(b_0)$.
For a general degree-one line bundle $F$, replacing $q^*F$ by a
single fiber is only a numerical equivalence unless a corresponding
holomorphic isomorphism has been chosen. Also, a general degree-one
bundle on a genus-two curve need not have nonzero $H^1$.
The choice just made avoids both qualifications.

For Koike's modified example, take again a genus-two curve and set
$F=K_B$. Its degree is two. Serre duality on the curve gives
\[
 H^1(B,K_B)\cong H^0(B,\OO_B)^*\cong\C
\]
because $B$ is compact and connected
\cite[Chapter III, \S7]{Hartshorne1977}. A nonzero element therefore
gives another nonsplit extension of the required form. In this case
\begin{equation}\label{con:degree-two}
 L_0=\OO_X(C)\otimes q^*K_B,\qquad C^2=-2,\qquad L_0^2=2.
\end{equation}
Consequently Corollary~\ref{con:application} applies to both the original
degree-one example and this degree-two example. Their numerical
invariants differ; they are not being identified as the same ruled
surface.

\subsection{The bounded-potential question}

The question in \cite[\S4]{Zhang2025} concerns the existence of a
positive current with bounded potentials and minimal singularities
in the distinguished class for either example. We spell out exactly
what the construction supplies. Fix a smooth metric $h_0$ on $L_0$
and set $\theta=\cOne(L_0,h_0)$. The global function $u$ constructed
in \eqref{met:global-potential} is bounded and satisfies
\[
                      T=\theta+\ddc u\ge0 .
\]
In a holomorphic frame, $T=\ddc(\varphi_0+u)$, and
$\varphi_0+u$ is bounded on each relatively compact chart.
For any other $v\in\PSH(X,\theta)$, compactness gives
$\sup_Xv<\infty$ unless $v\equiv-\infty$, in which case the
comparison below is automatic. Since $u$ is bounded below,
\[
                v\le u+\bigl(\sup_Xv-\inf_Xu\bigr).
\]
Thus $T$ has minimal singularities by the defining comparison.
These are exactly the two requested properties. Moreover its
potentials are locally Lipschitz, a stronger continuity assertion.
The construction provides an affirmative answer for the data of
\eqref{con:degree-one} and \eqref{con:degree-two}.

This conclusion is an implication of the proofs in this manuscript;
it is not attributed to Grauert, Koike, or Zhang. In particular,
Zhang's December 2025 version is cited as posing the question, not
as already proving the bounded-potential statement. No assertion
about publication priority of the present elementary construction
is needed for this implication.

The assertion concerns existence of a representative with this
regularity. Any two potentials with minimal singularities differ
by a bounded function, by applying the defining comparison in both
directions. Bounded difference alone does not imply that every
minimal potential is continuous or Lipschitz; no such conclusion
is made here.

\subsection{Why the contracted surface is not K\"ahler}

The analytic contraction from Proposition~\ref{geo:contraction} is compatible
with the existence of the Lipschitz current on $X$. In particular,
one should not infer that the current supplies a K\"ahler form on
the singular surface $Y$. The following argument establishes the
opposite conclusion directly.

\begin{proposition}\label{con:nonkahler}
The normal compact surface $Y$ obtained by contracting $C$ is not
K\"ahler in the usual sense for complex spaces: it has no K\"ahler
form locally given by smooth strictly plurisubharmonic potentials
in ambient holomorphic embeddings.
\end{proposition}

\begin{proof}
Suppose such a K\"ahler form $\omega_Y$ exists. Pull it back along
$\pi:X\to Y$. The result is a smooth real closed semipositive
$(1,1)$-form $\beta$ on $X$. To justify smoothness even on $C$,
embed a neighborhood in $Y$ into a domain in $\C^N$ and represent
$\omega_Y$ there by $\ddc\Phi$, with $\Phi$ a smooth strictly
plurisubharmonic ambient function. Its pullback is
$\ddc(\Phi\circ\pi)$, where the composition uses the holomorphic
map into that ambient domain. This is a smooth form on $X$.
These local pullbacks agree, and they are semipositive by the
chain rule for the complex Hessian. They are strictly positive
away from $C$, since $\pi$ is biholomorphic there. Along $TC$,
the form vanishes because $\pi|_C$ is constant. Therefore
\begin{equation}\label{con:beta-intersections}
                    \int_C\beta=0,\qquad \int_f\beta>0
\end{equation}
for every fiber $f$ of $q$. For the second inequality, $f\setminus C$
is nonempty and $\beta$ is strictly positive on its tangent line;
its restriction to all of $f$ is nonnegative.

We next identify its real cohomology class. The needed instance
of the Leray--Hirsch theorem says that, for a $\P^1$-bundle and
a class restricting to the generator of $H^2(\P^1,\R)$,
the cohomology of the total space is a free module over that of
the base with basis that class and $1$
\cite[Theorem 5.11]{BottTu1982}. Here $\cOne(L_0)$ restricts to
that generator because $L_0|_f=\OO_{\P^1}(1)$.
Choose $\tau\in H^2(B,\R)$ with $\int_B\tau=1$. In degree two
the theorem gives the direct-sum description
\[
 H^2(X,\R)=
        \R\,\cOne(L_0)\ \oplus\ \R\,q^*\tau .
\]
Hence $[\beta]=a\cOne(L_0)+bq^*\tau$ for real numbers $a,b$.
Pairing with $C$ gives $b=0$, since
$L_0\cdot C=0$ and $\int_Cq^*\tau=1$. Pairing with $f$ then
gives $a=\int_f\beta>0$. Thus
$[\beta/a]=\cOne(L_0)$.

The smooth projective surface $X$ is compact K\"ahler.
The $\partial\bar\partial$-lemma on a compact K\"ahler manifold
says that a real closed $(1,1)$-form with zero real de Rham
class is $\ddc$ of a smooth real function
\cite[Chapter VI, Lemma 8.6]{DemaillyCADG}.
Apply it to $\beta/a-\theta$, where $\theta$ is the curvature
of any smooth Hermitian metric $h_0$ on $L_0$. There is a smooth
real function $v$ such that
\[
                    \beta/a-\theta=\ddc v .
\]
The metric $h_0e^{-v}$ on $L_0$ is smooth and has semipositive
curvature $\beta/a$. This contradicts
Theorem~\ref{obs:c1-obstruction}. Hence $\omega_Y$ cannot exist.
\end{proof}

\begin{remark}\label{con:regularity-meaning}
The metric constructed here has near $C$ a weight of the form
$\varphi_0+M|w|e^{\lambda/2}$. Its failure to be $C^1$ is visible
on a real normal line $w=t$: the two one-sided derivatives of its
cusp term at $t=0$ are $Me^{\lambda(z,0)/2}$ and
$-Me^{\lambda(z,0)/2}$. Adding the smooth function $\varphi_0$
does not remove this jump. The obstruction theorem says much
more than this observation about our formula: no different
semipositive metric can have $C^1$ weights near $C$.
Conversely, the non-K\"ahlerness of $Y$ does not contradict the
positive current on $X$: the latter is not a smooth K\"ahler
form near $C$, and the construction gives no smooth strictly
plurisubharmonic potentials descending to the contracted point.
\end{remark}

% End consequences.tex
% Begin history.tex
\section{Historical context and the scope of the conclusion}
\label{hist:context}

The geometric construction, the obstruction to smooth semipositivity, and the
bounded-potential question arose at different times.  We distinguish these
statements, and also distinguish Grauert's degree-one example from Koike's
degree-two modification.  In the notation used here,
\[
  0\longrightarrow F\longrightarrow E\longrightarrow\mathcal O_B
     \longrightarrow0,\qquad
  X=\mathbb P_{\mathrm{lines}}(E),\qquad C=\mathbb P_{\mathrm{lines}}(F),
  \qquad L_0=\mathcal O_X(C)\otimes q^*F,
\]
where $q:X\to B$ is the projection, $\deg F=d>0$, and the extension class
$\xi\in H^1(B,F)$ is nonzero.  These conventions identify the normal bundle
as $N_{C/X}=F^{-1}$ and give $C^2=-d$.

\subsection{Grauert's construction and the degree-one choice}
\label{hist:grauert}

Grauert's example occurs in \cite[\S4, paragraph~8(d), pp.~365--366]{Grauert1962}.
Its starting data are a genus-two curve $B$, a degree-one line bundle $F$, and
a nontrivial affine-line bundle modeled on $F$.  The affine bundle is an
$F$-torsor: differences of local sections are sections of $F$, but its
transition functions also contain translations.  Compactifying each affine
fiber by one point gives the ruled surface above; the added section is $C$.
Thus the affine bundle itself is not a locally free sheaf, whereas $E$ is a
rank-two locally free sheaf.

The cohomological hypothesis $H^1(B,F)\ne0$ is essential.  It does not follow
merely from $\deg F=1$.  One explicit valid choice is
$F=\mathcal O_B(b)$ for a point $b\in B$: here $h^0(B,F)=1$, and
Riemann--Roch gives $h^1(B,F)=1$.  For comparison, a degree-one bundle outside
the effective locus has both groups zero.  The qualification concerning
$H^1(B,F)$ is also recorded in note~T28 of the supplied English translation
of Grauert's paper.  Choosing $\mathcal O_B(b)$ resolves it without altering
the intended example.

Since $C^2=-1$, Grauert's contraction theorem produces a proper holomorphic
map from $X$ to a normal compact complex surface, contracting precisely $C$
to one point.  The relevant contraction results are proved in
\cite[\S3 and \S4, paragraph~8(e)]{Grauert1962}; a modern surface treatment is
\cite[Chapter~III, \S2]{BHPV2004}.  The contracted point is singular, despite
$C^2=-1$, because $C$ has genus two.  Grauert proves that the target is not
projective: a hyperplane section avoiding the contracted point would lift
to a compact multisection of the affine bundle, and the fiberwise average
of its points would give a section of the nontrivial torsor.  In particular,
Grauert's original argument concerns analytic contraction and
nonprojectivity; it does not establish an obstruction to continuous
semipositive metrics.

\subsection{The questions of Filip--Tosatti and the obstruction of Koike}
\label{hist:ft-koike}

Filip--Tosatti revisit this surface in
\cite[\S2, especially Problem~2.2, p.~2989]{FilipTosatti2018}.  For an ample
line bundle $A$ they consider, in additive divisor notation,
\[
  A-\frac{A\cdot C}{C^2}\,C.
\]
Its intersection with $C$ is zero, while the class is nef and big.  Their
question asks for the failure of \emph{Hermitian semipositivity}, meaning the
existence of a smooth Hermitian metric with semipositive curvature.  When
$C^2=-1$ this displayed expression is integral.  If the construction is
written with $C^2=-d$ for general $d$, it is initially a rational divisor
class; denominators must be cleared before it is treated as a line bundle.
This issue does not arise for the explicitly defined integral bundle $L_0$.
Filip--Tosatti also explain that the contracted surface is non-K\"ahler,
referring to \cite{Moishezon1975}.  That assertion is not needed for the
metric construction in the present manuscript.

Koike develops first-order obstructions to linearizing a line bundle along
a compact submanifold in \cite{Koike2021}.  His Theorem~1.4 shows that a
nonzero class
\[
  u_1(C,X,L)\in H^1(C,N_{C/X}^{-1})
\]
precludes smooth semipositivity when $C$ is compact K\"ahler and $L|_C$ is
topologically trivial.  This extends an obstruction-theoretic viewpoint
originating in Ueda's study of curves with topologically trivial normal
bundle \cite{Ueda1983}.  The normal bundle of the present curve is negative,
so Ueda's original hypotheses are not themselves the hypotheses of the
application here.

In \cite[Example~4.1, pp.~2262--2263]{Koike2021}, Koike takes a genus-two curve
and $F=K_B$.  Thus his explicit model has $d=2$, $C^2=-2$, and
$L_0^2=2$.  A cocycle obtained from the hyperbolic metric represents a
nonzero class in $H^1(B,K_B)$, and the first-order obstruction of $L_0$ is
nonzero.  This proves that this modified Grauert example is nef and big
but has no smooth semipositive metric.  It is important not to identify
its degree-two data with Grauert's original degree-one data.

\subsection{Zhang's formulation and the conclusion proved here}
\label{hist:zhang}

In \cite{Zhang2025}, Zhang returns to the degree-one construction and to
the precise line bundle $L_0=q^*F\otimes\mathcal O_X(C)$.  The calculation
in \S3 identifies its first obstruction with the nontrivial affine-bundle
class and applies Koike's theorem to exclude a smooth semipositive metric.
Section~4 asks whether the first Chern class on the original or modified
Grauert example admits a positive current with minimal singularities and
bounded local potentials.  The version cited here is arXiv:2512.23446v1,
submitted on 29 December 2025; this formulation is attributed to that
specific version.

Theorem~\ref{main:regularity} and the global construction in
Theorem~\ref{met:lipschitz-global} give an affirmative answer for the data treated
here: $L_0$ has a positive-curvature Hermitian metric with locally Lipschitz
weights.  This includes both the original degree-one case and the
degree-two canonical-bundle case, and more generally every positive-degree
$F$ and nonzero extension class satisfying the stated hypotheses.
Theorem~\ref{obs:c1-obstruction} also excludes a metric with
$C^1$ weights and semipositive curvature.  These regularity conclusions
are deductions proved in this manuscript; they are not statements being
attributed to Grauert, Filip--Tosatti, Koike, or Zhang.  No claim of
bibliographic priority is intended.

Boundedness implies minimal singularities in this situation for an
elementary reason.  Fix a smooth reference metric and write a
positive-curvature metric as a global quasi-plurisubharmonic weight
$\varphi$.  Every competing such weight $\psi$ is bounded above on the
compact surface.  If $\varphi$ is bounded below, then
$\psi\leq\varphi+O(1)$, which is the defining comparison for minimal
singularities.  This terminology and its general framework are developed
in \cite[\S1]{DPS2001} and \cite{BEGZ2010}.  Thus the bounded-potential
conclusion concerns the same metric and class, without a separate
existence theorem for minimal singularities.

\subsection{What the standard tools do and do not supply}
\label{hist:tools}

The projectivity of $X$ follows by algebraizing the extension on the
projective curve using GAGA \cite{Serre1956}.  The projective-bundle and
ruled-surface calculations are standard
\cite[Chapters~II and~V]{Hartshorne1977}; here the convention of
parametrizing lines is fixed explicitly.  Ampleness of the auxiliary bundle
used in the construction is checked by the surface numerical criterion
\cite{Nakai1963,Lazarsfeld2004}.  These algebraic facts provide a smooth
positive metric away from the exceptional-curve difficulty; they do not
give smooth semipositivity of $L_0$.

The analytic patching uses maxima of plurisubharmonic functions and
removable singularities for locally bounded-above plurisubharmonic
functions.  Suitable precise versions are
\cite[Chapter~I, Lemma~5.17 and Theorem~5.24]{DemaillyCADG}; regularized
maxima and their preservation of Levi lower bounds appear in
\cite[Chapter~I, Lemma~5.18]{DemaillyCADG}, within the regularization theory
of \cite{Richberg1968}.  The decisive local estimate and the $C^1$
obstruction are supplied in the present proofs.  In particular, these
regularization results do not authorize smoothing the final metric
across $C$ while preserving its semipositive curvature.

There is no general implication from nefness to bounded semipositive
weights.  The nonsplit extension over an elliptic curve in
\cite[Example~1.7]{DPS1994} gives a nef, non-big tautological bundle;
\cite[Corollary~4.1 and Example~4.2]{Koike2015} explains its unbounded
minimal-metric behavior and a related nef and big threefold example.
Koike's positive results for negative curves in that paper impose, among
other conditions, a holomorphic tubular-neighborhood hypothesis.  Such a
hypothesis is not assumed here and cannot be silently substituted for the
nontrivial affine transitions.

% End history.tex

\begin{thebibliography}{BHPV04}

\bibitem[BEGZ10]{BEGZ2010}
S{\'e}bastien Boucksom, Philippe Eyssidieux, Vincent Guedj, and Ahmed Zeriahi.
\newblock {Monge--Amp\`ere} equations in big cohomology classes.
\newblock {\em Acta Mathematica}, 205(2):199--262, 2010.
\newblock \href{https://doi.org/10.1007/s11511-010-0054-7}{doi:\nolinkurl{10.1007/s11511-010-0054-7}}.

\bibitem[BHPV04]{BHPV2004}
Wolf~P. Barth, Klaus Hulek, Chris A.~M. Peters, and Antonius {Van de Ven}.
\newblock {\em Compact Complex Surfaces}, volume~4 of {\em Ergebnisse der Mathematik und ihrer Grenzgebiete, 3. Folge}.
\newblock Springer-Verlag, Berlin, second enlarged edition, 2004.
\newblock \href{https://doi.org/10.1007/978-3-642-57739-0}{doi:\nolinkurl{10.1007/978-3-642-57739-0}}.

\bibitem[BT82]{BottTu1982}
Raoul Bott and Loring~W. Tu.
\newblock {\em Differential Forms in Algebraic Topology}, volume~82 of {\em Graduate Texts in Mathematics}.
\newblock Springer-Verlag, New York, 1982.
\newblock \href{https://doi.org/10.1007/978-1-4757-3951-0}{doi:\nolinkurl{10.1007/978-1-4757-3951-0}}.

\bibitem[Dem12a]{Demailly2012}
Jean-Pierre Demailly.
\newblock {\em Analytic Methods in Algebraic Geometry}, volume~1 of {\em Surveys of Modern Mathematics}.
\newblock International Press and Higher Education Press, Somerville, MA and Beijing, 2012.
\newblock \href{https://intlpress.com/site/pub/files/preview/bookpubs/00000371.pdf}{Publisher's preview}.

\bibitem[Dem12b]{DemaillyCADG}
Jean-Pierre Demailly.
\newblock {\em Complex Analytic and Differential Geometry}.
\newblock Institut Fourier, Universit{\'e} de Grenoble, 2012.
\newblock Online monograph, version of 21 June 2012; \href{https://www-fourier.univ-grenoble-alpes.fr/~demailly/manuscripts/agbook.pdf}{author's PDF}.

\bibitem[DPS94]{DPS1994}
Jean-Pierre Demailly, Thomas Peternell, and Michael Schneider.
\newblock Compact complex manifolds with numerically effective tangent bundles.
\newblock {\em Journal of Algebraic Geometry}, 3(2):295--345, 1994.
\newblock \href{https://eref.uni-bayreuth.de/id/eprint/59615/}{Institutional record}.

\bibitem[DPS01]{DPS2001}
Jean-Pierre Demailly, Thomas Peternell, and Michael Schneider.
\newblock Pseudo-effective line bundles on compact {K\"ahler} manifolds.
\newblock {\em International Journal of Mathematics}, 12(6):689--741, 2001.
\newblock \href{https://doi.org/10.1142/S0129167X01000861}{doi:\nolinkurl{10.1142/S0129167X01000861}}.

\bibitem[FT18]{FilipTosatti2018}
Simion Filip and Valentino Tosatti.
\newblock Smooth and rough positive currents.
\newblock {\em Annales de l'Institut Fourier}, 68(7):2981--2999, 2018.
\newblock \href{https://doi.org/10.5802/aif.3234}{doi:\nolinkurl{10.5802/aif.3234}}.

\bibitem[Gra62]{Grauert1962}
Hans Grauert.
\newblock {{\"U}ber Modifikationen und exzeptionelle analytische Mengen}.
\newblock {\em Mathematische Annalen}, 146(4):331--368, 1962.
\newblock \href{https://doi.org/10.1007/BF01441136}{doi:\nolinkurl{10.1007/BF01441136}}.

\bibitem[Har77]{Hartshorne1977}
Robin Hartshorne.
\newblock {\em Algebraic Geometry}, volume~52 of {\em Graduate Texts in Mathematics}.
\newblock Springer-Verlag, New York--Heidelberg, 1977.
\newblock \href{https://doi.org/10.1007/978-1-4757-3849-0}{doi:\nolinkurl{10.1007/978-1-4757-3849-0}}.

\bibitem[Koi15]{Koike2015}
Takayuki Koike.
\newblock On minimal singular metrics of certain class of line bundles whose section ring is not finitely generated.
\newblock {\em Annales de l'Institut Fourier}, 65(5):1953--1967, 2015.
\newblock \href{https://doi.org/10.5802/aif.2978}{doi:\nolinkurl{10.5802/aif.2978}}.

\bibitem[Koi21]{Koike2021}
Takayuki Koike.
\newblock Linearization of transition functions of a semi-positive line bundle along a certain submanifold.
\newblock {\em Annales de l'Institut Fourier}, 71(5):2237--2271, 2021.
\newblock \href{https://doi.org/10.5802/aif.3439}{doi:\nolinkurl{10.5802/aif.3439}}.

\bibitem[Laz04]{Lazarsfeld2004}
Robert Lazarsfeld.
\newblock {\em Positivity in Algebraic Geometry {I}: Classical Setting: Line Bundles and Linear Series}, volume~48 of {\em Ergebnisse der Mathematik und ihrer Grenzgebiete, 3. Folge}.
\newblock Springer-Verlag, Berlin, 2004.
\newblock \href{https://doi.org/10.1007/978-3-642-18808-4}{doi:\nolinkurl{10.1007/978-3-642-18808-4}}.

\bibitem[Moi75]{Moishezon1975}
Boris~G. Moishezon.
\newblock Singular {K\"ahlerian} spaces.
\newblock In {\em Manifolds---Tokyo 1973: Proceedings of the International Conference on Manifolds and Related Topics in Topology, Tokyo, 1973}, pages 343--351. University of Tokyo Press, Tokyo, 1975.

\bibitem[Nak63]{Nakai1963}
Yoshikazu Nakai.
\newblock A criterion of an ample sheaf on a projective scheme.
\newblock {\em American Journal of Mathematics}, 85(1):14--26, 1963.
\newblock \href{https://doi.org/10.2307/2373180}{doi:\nolinkurl{10.2307/2373180}}.

\bibitem[Ric68]{Richberg1968}
Rolf Richberg.
\newblock {Stetige streng pseudokonvexe Funktionen}.
\newblock {\em Mathematische Annalen}, 175:257--286, 1968.
\newblock \href{https://doi.org/10.1007/BF02063212}{doi:\nolinkurl{10.1007/BF02063212}}.

\bibitem[Ser56]{Serre1956}
Jean-Pierre Serre.
\newblock G{\'e}om{\'e}trie alg{\'e}brique et g{\'e}om{\'e}trie analytique.
\newblock {\em Annales de l'Institut Fourier}, 6:1--42, 1956.
\newblock \href{https://doi.org/10.5802/aif.59}{doi:\nolinkurl{10.5802/aif.59}}.

\bibitem[Ued83]{Ueda1983}
Tetsuo Ueda.
\newblock On the neighborhood of a compact complex curve with topologically trivial normal bundle.
\newblock {\em Journal of Mathematics of Kyoto University}, 22(4):583--607, 1983.
\newblock \href{https://doi.org/10.1215/kjm/1250521670}{doi:\nolinkurl{10.1215/kjm/1250521670}}.

\bibitem[Zha25]{Zhang2025}
Yangyang Zhang.
\newblock On the non-semipositivity of a nef and big line bundle on {Grauert's} example, 2025.
\newblock \href{https://arxiv.org/abs/2512.23446v1}{arXiv:2512.23446v1}, submitted 29 December 2025.

\end{thebibliography}

\end{document}
